% !TeX program = xelatex
% !TeX encoding = UTF-8
\documentclass[UTF8,11pt,openany]{ctexbook}
\usepackage[a4paper,left=27mm,right=25mm,top=27mm,bottom=27mm]{geometry}
\usepackage{amsmath,amssymb,amsthm,mathtools,bm}
\usepackage{booktabs,array,longtable,enumitem}
\usepackage{graphicx,xcolor}
\usepackage[most]{tcolorbox}
\usepackage{hyperref}
\usepackage{fancyhdr}
\usepackage{listings}
\usepackage{needspace}
\usepackage{placeins}
\usepackage{tikz}
\usepackage{pgfplots}
\pgfplotsset{compat=1.18}
\usetikzlibrary{positioning}
\hypersetup{colorlinks=true,linkcolor=blue!45!black,urlcolor=blue!50!black,pdftitle={从洛伦兹协变到双曲守恒律：Maxwell 方程的几何与计算},pdfauthor={Harvey Tan}}
\setlength{\parindent}{2em}
\setlength{\parskip}{0.18em}
\setlist[itemize]{itemsep=0.15em,topsep=0.25em}
\setlist[enumerate]{itemsep=0.15em,topsep=0.25em}
\pagestyle{fancy}
\fancyhf{}
\fancyhead[LE,RO]{\small\nouppercase{\leftmark}}
\fancyfoot[C]{\thepage}
\setlength{\headheight}{15pt}
\renewcommand{\headrulewidth}{0.4pt}
\numberwithin{equation}{chapter}
\newtheorem{theorem}{定理}[chapter]
\newtheorem{proposition}[theorem]{命题}
\newtheorem{lemma}[theorem]{引理}
\theoremstyle{definition}
\newtheorem{definition}[theorem]{定义}
\newtheorem{example}[theorem]{例}
\theoremstyle{remark}
\newtheorem{remark}[theorem]{说明}
\newcommand{\dd}{\mathrm d}
\newcommand{\RR}{\mathbb R}
\newcommand{\Om}{\Omega}
\newcommand{\vol}{\mathrm{vol}}
\newcommand{\grad}{\operatorname{grad}}
\newcommand{\Div}{\operatorname{div}}
\newcommand{\curl}{\operatorname{curl}}
\newcommand{\Tr}{\operatorname{tr}}
\newcommand{\starS}{\star_3}
\newcommand{\norm}[1]{\left\lVert #1\right\rVert}
\newcommand{\abs}[1]{\left\lvert #1\right\rvert}
\newcommand{\ip}[2]{\left\langle #1,#2\right\rangle}
\newcommand{\mat}[1]{\begin{pmatrix}#1\end{pmatrix}}
\definecolor{inkblue}{RGB}{35,65,108}
\lstset{basicstyle=\ttfamily\small,frame=single,breaklines=true,columns=fullflexible,keepspaces=true,upquote=true}
\newtcolorbox{chaptersummary}{colback=blue!3,colframe=blue!45!black,title={本章总结},fonttitle=\bfseries,boxrule=0.6pt,arc=2mm}
\newtcolorbox{studentinput}{colback=red!2,colframe=red!55!black,title={这里需要你自己补充},fonttitle=\bfseries,boxrule=0.7pt,arc=2mm}
\newtcolorbox{askai}{colback=violet!3,colframe=violet!55!black,title={你应该问 AI},fonttitle=\bfseries,boxrule=0.7pt,arc=2mm}
\newtcolorbox{checkpoint}{colback=green!3,colframe=green!45!black,title={理解检查点},fonttitle=\bfseries,boxrule=0.6pt,arc=2mm}
\newtcolorbox{engineernote}{colback=gray!5,colframe=gray!60!black,title={工程师视角},fonttitle=\bfseries,boxrule=0.6pt,arc=2mm}

\newcommand{\studyroute}[3]{%
  \par\smallskip\noindent\textbf{展开与检索。}#1\par
  \noindent\textbf{关键词：}#2。\par
  \noindent\textbf{AI 展开：}#3\par
}
\raggedbottom

\title{从洛伦兹协变到双曲守恒律\\[0.6em]\large Maxwell 方程的几何结构、初值问题与数值求解}
\author{Harvey Tan\\[0.6em]\large 计算物理学简明教程\quad 第 2 卷}
\date{2026 年 11 月}

\begin{document}
\frontmatter
\maketitle
\chapter*{前言}
\addcontentsline{toc}{chapter}{前言}
\begin{quote}
\itshape
此书献给朗道与栗夫席次。\\
感谢你们为无数后来者引路，让我们得见理论之美，并为之震撼。
\end{quote}
\begin{tcolorbox}[colback=red!3,colframe=red!65!black,title={使用要求：AI 讲解与代码实验},fonttitle=\bfseries,boxrule=0.8pt,arc=1mm,breakable]
\textbf{请勿仅凭正文孤立自学；本书须结合 AI 讲解与代码实验使用。}
正文提供高度压缩的知识地图。遇到难以理解的概念联系、推导或先备知识，
请携带当前符号、假设和具体疑问，让 AI 逐步展开；
同时运行配套代码或构造最小算例，改变参数并检查误差、守恒量或稳定性。
AI 的解释、推导与代码须对照正文假设、参考文献及数值结果核验。
\end{tcolorbox}
\noindent 第一卷建立了粒子与连续介质模型、偏微分方程和离散求解之间的联系。本卷以 Maxwell 方程为主线，研究场方程的几何结构及其对数值计算的影响。

本卷依次讨论：电磁场的四维表示；演化方程与约束的时空分解；对称双曲性、能量估计和守恒型；以及散度约束在数值离散中的保持。

本卷假定读者熟悉多元微积分、线性代数、常微分方程、基本偏微分方程及向量分析。第 0 章介绍后文所需的微分形式、函数空间与辛结构；第 5.5 章讨论时间推进，第 7 章从弱形式进入 Galerkin 与 DG。

本卷以三维 Maxwell 方程为物理输入，将其组织为四维几何表达。Lorentz 协变性约束方程的变换形式，而动力学内容需要由物理规律给定。

\paragraph{贯穿全文的约定} 除专门指出外，取真空、平直时空、无量纲化单位 $c=\varepsilon_0=\mu_0=1$；时空度规号差为 $(-,+,+,+)$；坐标为 $(t,x,y,z)$；四维取向为 $\dd t\wedge\dd x\wedge\dd y\wedge\dd z$；三维空间取向为 $\dd x\wedge\dd y\wedge\dd z$。空间向量与一形式之间的升降指标使用欧氏度规。所有经典推导先假设足够光滑；弱解和离散问题在后文讨论。符号约定一经改变，若干 Hodge 星与分量公式会出现符号差异，但物理方程不变。

\tableofcontents

\mainmatter
\setcounter{chapter}{-1}
\chapter{数学准备：从向量分析走向几何与 PDE}
\label{ch:prelim}
\begin{chaptersummary}
第一卷介绍了连续方程与离散求解，本章补充研究 Maxwell 结构所需的语言：微分形式组织微分算子，函数空间与能量估计描述初值问题，辛结构用于讨论长期时间推进。这些工具将在后续章节分别连接约束传播、弱形式和离散演化。核心是区分方程的几何结构、解的正则性和数值演化结构。
\end{chaptersummary}
\studyroute{这里略去外代数、迹定理与辛几何的一般证明，保留后文使用的定义和算子关系。}{differential forms; Hodge star; Sobolev traces; symplectic map}{把 $\dd^2=0$ 翻译成向量恒等式，再说明哪些步骤使用度规、哪些仅使用微分结构。}
\section{数学对象与分析目标}
数学物理方法通常教我们求一个给定 PDE 的解；现代 PDE 理论还关心：什么是合适的初值？解是否存在、唯一、稳定？方程中有没有必须始终满足的约束？数值离散是否保留其结构？为讨论这些问题，我们依次补充\emph{微分形式}、\emph{时空分解}、\emph{函数空间}和\emph{守恒律}的基础语言。

\begin{center}
\begin{tikzpicture}[>=stealth, node distance=8mm]
\node[draw,rounded corners,align=center,minimum width=34mm] (v) {向量分析\\$\grad,\curl,\Div$};
\node[draw,rounded corners,align=center,right=of v,minimum width=34mm] (f) {微分形式\\$\dd,\wedge,\star$};
\node[draw,rounded corners,align=center,right=of f,minimum width=34mm] (p) {初值问题\\双曲性与能量};
\draw[->] (v)--(f);\draw[->] (f)--(p);
\end{tikzpicture}
\end{center}

\section{多重线性代数：为何需要反对称对象}
在三维，向量叉乘非常便利，但它依赖维数、度量和取向；在四维，反对称张量更加自然。一个 $k$ 形式可以理解为一个对 $k$ 个向量输入呈线性、且交换两个输入就变号的函数。

一形式的例子为
\begin{equation}
\alpha=a_x\,\dd x+a_y\,\dd y+a_z\,\dd z.
\end{equation}
它作用于向量 $v=v^i\partial_i$ 得到 $\alpha(v)=a_iv^i$。二形式的例子为
\begin{equation}
\beta=b_x\,\dd y\wedge\dd z+b_y\,\dd z\wedge\dd x+b_z\,\dd x\wedge\dd y.
\end{equation}
$\wedge$ 称为外积，满足 $\dd x\wedge\dd y=-\dd y\wedge\dd x$，尤其 $\dd x\wedge\dd x=0$。一般地，$\alpha\in\Omega^p$、$\beta\in\Omega^q$ 时，
\begin{equation}
\alpha\wedge\beta=(-1)^{pq}\beta\wedge\alpha,\qquad
\alpha\wedge\beta\in\Omega^{p+q}.
\end{equation}
一个 $n$ 维空间中独立 $k$ 形式的数目是 $\binom nk$；比如四维有 6 个独立二形式和 4 个独立三形式。$\Omega^k$ 表示所有光滑 $k$ 形式构成的集合（更准确地说是向量空间）。

\section{外微分：把 grad、curl、div 合成一个算子}
定义外微分 $\dd:\Omega^k\to\Omega^{k+1}$。在欧氏三维空间，
\begin{align}
\dd f&=\partial_x f\,\dd x+\partial_yf\,\dd y+\partial_zf\,\dd z,\\
\dd(a_i\dd x^i)&=(\curl\bm a)_x\,\dd y\wedge\dd z+(\curl\bm a)_y\,\dd z\wedge\dd x
 +(\curl\bm a)_z\,\dd x\wedge\dd y,\\
\dd(b_x\dd y\wedge\dd z+\cdots)&=(\Div\bm b)\,\dd x\wedge\dd y\wedge\dd z.
\end{align}
这里等号经过了下节的欧氏识别：一形式对应通常的向量分量，二形式对应穿过有向面积的通量分量。最关键的恒等式是
\begin{equation}\label{eq:d2zero}
\boxed{\dd^2=0.}
\end{equation}
在熟悉的向量记号中，这同时包含 $\curl\grad f=0$、$\Div\curl\bm a=0$。外微分还满足分次 Leibniz 法则：
\begin{equation}
\dd(\alpha\wedge\beta)=\dd\alpha\wedge\beta+(-1)^p\alpha\wedge\dd\beta,
\quad \alpha\in\Omega^p.
\end{equation}
\begin{remark}
$\dd^2=0$ 可由混合偏导可交换加上反对称性推出。对于更一般的几何对象，该恒等式并不需要先定义长度和夹角。
\end{remark}

\begin{studentinput}
请解释 $\curl\grad f=0$ 和 $\Div\curl\bm a=0$ 如何同时来自 $\dd^2=0$。为什么这一关系与特定 Maxwell 解无关？若离散算子满足 $D_hC_h\ne0$，磁散度可能怎样变化？

\textbf{你的答案：}\vspace{2cm}
\end{studentinput}

\section{Hodge 星：哪些步骤需要度规}
\begin{askai}
请从三维欧氏空间的 $\dd x,\dd y,\dd z$ 出发，逐步计算 $\starS\dd x$ 和 $\starS(\dd y\wedge\dd z)$，并指出取向与度规各在哪一步发挥作用。

进一步展开：将度规缩放后，外微分 $\dd$ 与 Hodge 星各会如何变化？请给出一个具体例子。
\end{askai}
外微分由光滑结构定义；借助 Hodge 星将一形式对应到二形式，还需要\emph{度规与取向}。在有度规和取向的 $n$ 维空间上，Hodge 星将 $k$ 形式映为 $(n-k)$ 形式：
\begin{equation}
\star:\Omega^k\to\Omega^{n-k},\qquad
\alpha\wedge\star\beta=\ip{\alpha}{\beta}\,\vol_n.
\end{equation}
三维欧氏情形有
\begin{equation}
\starS\dd x=\dd y\wedge\dd z,\quad
\starS\dd y=\dd z\wedge\dd x,\quad
\starS\dd z=\dd x\wedge\dd y,
\end{equation}
且 $\starS^2=1$。因而 $\bm E$ 的一形式 $E=E_i\dd x^i$ 可通过 $\starS E$ 变成面积通量二形式。$\dd$ 与度规无关，$\star$ 则由度规和取向定义。在 Lorentz 时空，后者还反映时间方向的负号。

\section{积分定理与拓扑的简短提醒}
广义 Stokes 定理把基本的微积分定理、Green 公式、Gauss 公式、Stokes 公式统一为
\begin{equation}\label{eq:stokes}
\int_\Omega\dd\alpha=\int_{\partial\Omega}\alpha.
\end{equation}
这也是守恒律积分形式的几何来源。对次数为正的微分形式，在可缩邻域内，$\dd\alpha=0$ 蕴含 $\alpha=\dd\gamma$（Poincar\'e 引理）；非平凡拓扑可能阻碍这一表示在全局成立。后文的 $F=\dd A$ 因而是局部陈述，全局表示需要相应拓扑条件。

\section{时间、空间与初值问题}
从 ODE 的 $\dot u=f(u)$ 推广到 PDE，典型演化问题写成
\begin{equation}
\partial_t u+\mathcal L_xu=f,\qquad u(0,x)=u_0(x).
\end{equation}
这里 $t$ 是演化时间，$x$ 是空间变量；这种选择称为\textbf{时空 $3+1$ 分解}或\emph{时间分叶}，与求解析解时的分离变量法不同。选定时间后，协变方程可分为含时间导数的演化式与不含时间导数的约束式。

\section{函数空间与能量估计}
\begin{studentinput}
设分片光滑的电场跨界面有跳跃。哪些切向或法向跳跃会分别影响其分布意义下的旋度和散度？结合后文说明为何要分别使用 $H(\curl)$ 与 $H(\Div)$。

\textbf{你的答案：}\vspace{2cm}
\end{studentinput}
函数空间以整个函数为元素，用范数描述其大小与正则性。对 $\Omega\subset\RR^3$，定义
\begin{equation}
\norm{u}_{L^2(\Omega)}^2=\int_\Omega |u(x)|^2\,\dd x,
\quad
\norm{u}_{H^1(\Omega)}^2=\norm{u}_{L^2}^2+\norm{\nabla u}_{L^2}^2.
\end{equation}
$L^2$ 对应“总平方可积”，$H^1$ 还要求一阶\emph{弱导数}平方可积。弱导数的想法是把导数转移到测试函数：如果
\begin{equation}
\int_\Omega u\,\partial_i\varphi\,\dd x=-\int_\Omega v_i\varphi\,\dd x
\quad\text{对所有 }\varphi\in C_c^\infty(\Omega),
\end{equation}
则称 $v_i$ 是 $u$ 的弱导数。这允许我们处理分片光滑甚至带跳跃的数值解。

Maxwell 方程通常还涉及
\begin{align}
H(\curl;\Omega)&=\{\bm u\in L^2(\Omega)^3:\curl\bm u\in L^2(\Omega)^3\},\\
H(\Div;\Omega)&=\{\bm u\in L^2(\Omega)^3:\Div\bm u\in L^2(\Omega)\},
\end{align}
其中旋度、散度按分布意义理解。\emph{函数空间确定场与微分算子存在的意义。}完备性、切向迹与法向迹定理的证明在此略去，可检索 Sobolev graph spaces、tangential trace 和 normal trace。

\begin{definition}[适定性的工作定义]
一个初值问题若对允许的初值有解、解唯一，且在选定范数下连续依赖初值（与源项），就称在该设置下适定。何谓“允许”，由边界条件和所选函数空间决定。
\end{definition}
后文的 $L^2$ 能量估计揭示真空 Maxwell 演化的稳定性结构；完整适定性还依赖初值正则性和边界条件。


\section{辛结构与 Hamilton 系统}\label{sec:pre-symp}
\paragraph{从位置与动量出发。}
对于自由度为 $m$ 的有限维系统，记广义坐标 $q\in\RR^m$、共轭动量
$p\in\RR^m$，合成 $z=(q,p)^\mathsf T\in\RR^{2m}$。给定 Hamilton 函数
$H(q,p)$，正则 Hamilton 方程为
\begin{equation}\label{eq:canonicalham}
\dot z=\mathbb J\nabla H(z),\qquad
\mathbb J=\begin{pmatrix}0&I_m\\-I_m&0\end{pmatrix},
\quad
\dot q=\partial_pH,\quad \dot p=-\partial_qH.
\end{equation}
熟悉的无阻尼振子 $H=(p^2+\omega^2q^2)/2$ 对应
$\dot q=p,\;\dot p=-\omega^2q$；轨迹是相平面上的椭圆。

\paragraph{什么是辛？}
正则二形式
\begin{equation}
\omega=\sum_{i=1}^m\dd q_i\wedge\dd p_i
\end{equation}
度量的是相空间中的\emph{有向面积元素}，而不是欧氏长度。光滑映射
$\Phi:z\mapsto\Phi(z)$ 称为\emph{保辛}，如果
\begin{equation}\label{eq:sympmap}
(D\Phi)^\mathsf T\mathbb J(D\Phi)=\mathbb J.
\end{equation}
正则 Hamilton 系统的精确时间流保持此结构。保辛数值积分器的设计目标，
就是让\emph{一步映射}也满足 \eqref{eq:sympmap}。
辛结构保持、能量精确守恒和约束保持是\textbf{三个不同的命题}；在一般非线性 Hamilton 系统中，
保辛并不意味着每步能量完全不变。

\paragraph{时间推进为何会影响长期计算？}
显式 Euler 作用于振子，
\begin{equation}
\begin{pmatrix}q^{n+1}\\p^{n+1}\end{pmatrix}
=\begin{pmatrix}1&\Delta t\\-\omega^2\Delta t&1\end{pmatrix}
\begin{pmatrix}q^n\\p^n\end{pmatrix},
\end{equation}
其矩阵行列式为 $1+\omega^2\Delta t^2>1$，导致相面积膨胀及非物理的
长期能量增长。辛 Euler 的一种写法为
\begin{equation}\label{eq:sympeuler}
 p^{n+1}=p^n-\Delta t\,\omega^2q^n,\qquad
 q^{n+1}=q^n+\Delta t\,p^{n+1}.
\end{equation}
它的一步映射保辛，在满足适当步长条件时不会产生上述 Euler 型的系统性能量膨胀。
后文将以隐式中点法与交错推进讨论 Maxwell 半离散系统。此处略去辛流形与后向误差分析；检索 symplectic map、modified Hamiltonian 可继续研究保辛性与长期能量行为的联系。

\section{练习}
\begin{enumerate}
\item 验证 $\dd(a\,\dd x+b\,\dd y)$ 的 $\dd x\wedge\dd y$ 系数为 $\partial_xb-\partial_ya$。
\item 从 $\dd^2f=0$ 恢复 $\curl\grad f=0$。
\item 解释为什么四维二形式有 6 个独立分量，但四维三形式只有 4 个独立分量。
\item 说明 $L^2$ 范数为何适合度量电磁场的能量。
\item 对振子 $H=(p^2+\omega^2q^2)/2$，验证显式 Euler 一步矩阵行列式大于 $1$，解释这与面积保持的关系。
\end{enumerate}

\chapter[从 Lorentz 协变性到四维 Maxwell 方程]{从 Lorentz 协变性\\到四维 Maxwell 方程}
\label{ch:lorentz}
\begin{chaptersummary}
三维 Maxwell 方程在不同惯性系中描述同一电磁现象，但 Lorentz 变换会混合时间与空间导数。本章先从一个分量的实际变换发现电场与磁场的混合，再用反对称张量统一六个分量，核对四维方程与电荷守恒。这一统一为下一章选择时间方向、区分演化与约束提供起点。
\end{chaptersummary}
\studyroute{这里给出沿坐标轴的 boost；略去任意方向的变换推导与作用量变分。Maxwell 方程作为物理输入。}{passive Lorentz boost; electromagnetic field transformation; covariant Maxwell equations}{取本章被动 boost，先令仅有 $E_y$ 非零，再令仅有 $B_z$ 非零；逐项计算变换后的场，并用 $-v$ 的逆变换核验。}

\paragraph{从一次惯性系变换开始}
同一电磁现象在不同惯性系中应满足相同的物理规律。我们的出发点是已经由实验建立的三维 Maxwell 方程：
\begin{equation}\label{eq:maxwell3}
\begin{aligned}
\partial_t\bm B+\curl\bm E&=0,&\Div\bm B&=0,\\
\partial_t\bm E-\curl\bm B&=-\bm j,&\Div\bm E&=\rho.
\end{aligned}
\end{equation}
设惯性系 $S'$ 相对 $S$ 沿 $+x$ 方向以速度 $v$ 运动，$|v|<1$。同一事件的坐标满足
\[
t'=\gamma(t-vx),\qquad x'=\gamma(x-vt),\qquad
y'=y,\quad z'=z,\qquad \gamma=(1-v^2)^{-1/2}.
\]
这是被动坐标变换 $x'=\Lambda x$；撇号表示同一物理场在新坐标中的分量，
未撇分量均在对应事件 $x=\Lambda^{-1}x'$ 取值，并非主动移动场分布。
链式法则给出
\[
\partial_t=\gamma(\partial_{t'}-v\partial_{x'}),\qquad
\partial_x=\gamma(\partial_{x'}-v\partial_{t'}).
\]
把它代入 Faraday 方程的 $z$ 分量，得到
\[
\begin{aligned}
0&=\partial_tB_z+\partial_xE_y-\partial_yE_x\\
 &=\partial_{t'}\!\left[\gamma(B_z-vE_y)\right]
  +\partial_{x'}\!\left[\gamma(E_y-vB_z)\right]-\partial_{y'}E_x.
\end{aligned}
\]
进入新时间导数的已经是 $B_z-vE_y$，而非单独的 $B_z$。
要把结果重新组织为新观察者的 Faraday 方程，就出现
$B_z'=\gamma(B_z-vE_y)$、$E_y'=\gamma(E_y-vB_z)$ 这两个组合。
其余分量及源方程还须同时满足；下面的张量构造将完成这一核对。

\paragraph{六个分量的共同变换}
将场分量排列为
\[
(F_{\mu\nu})=
\begin{pmatrix}
0&-E_x&-E_y&-E_z\\
E_x&0&B_z&-B_y\\
E_y&-B_z&0&B_x\\
E_z&B_y&-B_x&0
\end{pmatrix},\qquad F_{\nu\mu}=-F_{\mu\nu}.
\]
这里始终采用 $F_{0i}=-E_i$、$F_{ij}=\epsilon_{ijk}B_k$，$\epsilon_{123}=1$。
反对称性把两个有序指标的交换变为符号改变，独立分量恰为
$\binom42=6$，与三个电场和三个磁场分量相符。
令这一对象按协变二阶张量变换：
\[
F'_{\mu\nu}(x')=K^\alpha{}_\mu K^\beta{}_\nu F_{\alpha\beta}(x),
\qquad K^\alpha{}_\mu=\frac{\partial x^\alpha}{\partial x'^\mu}
=(\Lambda^{-1})^\alpha{}_\mu.
\]
逆 Jacobian 的时间--空间块满足
$K^0{}_0=K^1{}_1=\gamma$、$K^0{}_1=K^1{}_0=\gamma v$；
例如
\[
F'_{02}=\gamma F_{02}+\gamma vF_{12}
=-\gamma(E_y-vB_z).
\]
同样计算六个独立分量，得到
\[
\boxed{\begin{aligned}
E_x'&=E_x,& B_x'&=B_x,\\
E_y'&=\gamma(E_y-vB_z),& B_y'&=\gamma(B_y+vE_z),\\
E_z'&=\gamma(E_z+vB_y),& B_z'&=\gamma(B_z-vE_y).
\end{aligned}}
\]
这与刚才的分量计算一致。电场与磁场是同一四维场相对于观察者的两部分；
四维统一的依据是它们共同的变换规律，而非给两个三维矢量添上时间坐标。

\paragraph{从分量方程到协变方程}
用 $\eta=\operatorname{diag}(-1,1,1,1)$ 升指标后，
$F^{0i}=E_i$、$F^{i0}=-E_i$，空间分量不变。
因此
\[
\partial_\mu F^{0\mu}=\Div\bm E,\qquad
\partial_\mu F^{i\mu}=(\curl\bm B)_i-\partial_tE_i.
\]
它们分别等于 $\rho$ 和 $j_i$；反对称导数的 $(123)$ 与 $(0ij)$ 分量
则分别给出无磁单极约束和 Faraday 方程。四条三维方程由此组织为
\[
\partial_{[\lambda}F_{\mu\nu]}=0,\qquad
\partial_\mu F^{\nu\mu}=j^\nu,\qquad j^\nu=(\rho,\bm j).
\]
其中方括号表示指标的完全反对称化。
电流按四矢量变换 $j'^\nu=\Lambda^\nu{}_\beta j^\beta$，特别有
$\rho'=\gamma(\rho-vj_x)$、$j_x'=\gamma(j_x-v\rho)$。
由 $\partial'_\mu=K^\alpha{}_\mu\partial_\alpha$ 和 Lorentz 变换保持度规，
\[
\partial'_\mu F'^{\nu\mu}
=\Lambda^\nu{}_\beta\,\partial_\alpha F^{\beta\alpha}
=j'^\nu;
\]
齐次方程同样按三个协变指标整体变换。结合既定取向，
这些方程的微分形式表达就是式 \eqref{eq:4maxwell} 的
$\dd F=0$、$\dd\star F=J$，其中 $F$ 与 $J$ 分别按
式 \eqref{eq:defF}、\eqref{eq:defJ} 定义。
Lorentz 协变性限定方程的变换形式，Maxwell 动力学的具体内容由物理规律给定。

计算机推进初值时还需选定时间方向：四维方程中哪些分量成为演化式，
哪些成为 Gauss 约束，以及这些约束如何在连续演化和离散计算中保持，
正是第二章及后续数值章节要研究的问题。

\section{物理输入与协变表达}
式 \eqref{eq:maxwell3} 提供动力学，Lorentz 变换检验不同惯性观察者的描述是否相容。
开篇的计算据此识别了统一的场变量；接下来把它写成二形式，
使演化、约束和电荷守恒能够由同一组几何运算组织。
场的变换与协变方程也可参照
\href{https://ocw.mit.edu/courses/8-033-introduction-to-relativity-and-spacetime-physics-fall-2024/mit8_033_f24_lec_full.pdf}{Hughes 讲义第 11--12 讲}；
比较分量时仍以本卷的指标与符号约定为准。

\section{电磁场强度二形式}
\begin{askai}
取本章被动 boost，令 $S$ 中仅有 $E_y=E_0$ 非零。用逆 Jacobian
显式计算 $K^{\mathsf T}FK$，核对 $E_y'$、$B_z'$，再用 $-v$ 恢复原场。

进一步展开：为何仅要求 Lorentz 协变性不能唯一确定 Maxwell 动力学？请比较两个协变方程的例子。
\end{askai}
在 $M=\RR^{1,3}$ 上令 $\eta=\operatorname{diag}(-1,1,1,1)$。定义空间一形式 $E=E_x\dd x+E_y\dd y+E_z\dd z$ 与空间二形式
\begin{equation}
B=B_x\dd y\wedge\dd z+B_y\dd z\wedge\dd x+B_z\dd x\wedge\dd y.
\end{equation}
把它们合并为
\begin{equation}\label{eq:defF}
\boxed{F=B+E\wedge\dd t=B-\dd t\wedge E\in\Omega^2(M).}
\end{equation}
开篇的反对称矩阵正是这一二形式的坐标分量：
\begin{equation}
F=\frac12F_{\mu\nu}\,\dd x^\mu\wedge\dd x^\nu,\qquad
F_{0i}=-E_i,\quad F_{ij}=\epsilon_{ijk}B_k.
\end{equation}
张量变换律保证同一几何对象的两种坐标展开相等：
\begin{equation}
F=\frac12F_{\mu\nu}\,\dd x^\mu\wedge\dd x^\nu
=\frac12F'_{\mu\nu}\,\dd x'^\mu\wedge\dd x'^\nu.
\end{equation}
因此观察者改变的是电磁分解与分量，所描述的 $F$ 保持为同一对象。

\section{四维表达与三种不同的“次数”}
取时空 Hodge 星 $\star$，将四维电流 $j^\mu=(\rho,\bm j)$ 写成三形式
\begin{equation}\label{eq:defJ}
J=\rho\,\vol_3-\dd t\wedge\starS j^\flat,
\qquad \vol_3=\dd x\wedge\dd y\wedge\dd z.
\end{equation}
这里 $j^\flat=j_i\dd x^i$ 仅表示\emph{空间}电流的一形式。Maxwell 方程等价于
\begin{equation}\label{eq:4maxwell}
\boxed{\dd F=0,\qquad \dd\star F=J.}
\end{equation}
在上述约定下，它们也可写成张量方程
\begin{equation}
\partial_{[\lambda}F_{\mu\nu]}=0,
\qquad \partial_\mu F^{\nu\mu}=j^\nu.
\end{equation}
不同教材若把 $F_{0i}$ 定义为 $+E_i$，第二式的指标顺序或右端符号会相应调整。

时空维数为 4，电磁场 $F$ 是 2 形式，Maxwell 方程的各项是 3 形式；下一节的守恒恒等式涉及 4 形式。

\section{外微分的平方为零与电荷守恒}
\begin{checkpoint}
从 $\dd\star F=J$ 出发，说明 $\dd J=0$ 为什么是源项的相容条件。给定不满足电荷连续性方程的源，原方程组还能否有光滑解？
\end{checkpoint}
对 $\dd\star F=J$ 再取一次外微分：
\begin{equation}\label{eq:dJ}
\dd J=\dd^2\star F=0.
\end{equation}
此时 $\dd J\in\Omega^4(M)$。把式 \eqref{eq:defJ} 代入可得
\begin{equation}
\dd J=(\partial_t\rho+\Div\bm j)\,
\dd t\wedge\vol_3,
\end{equation}
即
\begin{equation}\label{eq:continuity}
\boxed{\partial_t\rho+\Div\bm j=0.}
\end{equation}
电荷守恒是四维 Maxwell 方程的\emph{相容性条件}。反过来，求带源 Maxwell 初值问题时，给定的 $\rho,\bm j$ 必须满足该条件。

\paragraph{局部势表示（选读）}
由于 $\dd F=0$，局部可令 $F=\dd A$，其中 $A=\bm A^\flat-\phi\dd t$。于是 $\bm B=\curl\bm A$、$\bm E=-\partial_t\bm A-\grad\phi$。规范变换 $A\mapsto A+\dd\chi$ 不改变 $F$。势表示连接作用量、电磁波方程与规范理论；后文继续采用场强变量。

\section{本章小结与练习}
\noindent\textbf{小结：}Lorentz 协变性说明为何寻找统一的四维电磁场；微分形式为这种统一提供恰当的语言；Maxwell 物理规律提供动力学内容。
\begin{enumerate}
\item 说明 $F$ 为什么恰有 6 个独立分量。
\item 用 $\dd^2=0$ 解释为什么 $J$ 不能任意给定。
\item 思考：如果 $\rho,\bm j$ 不满足连续性方程，$\dd\star F=J$ 是否有光滑解？
\end{enumerate}

\chapter{时空 \texorpdfstring{$3+1$}{3+1} 分解：演化与约束怎样出现}
\label{ch:31}
\begin{chaptersummary}
四维方程统一了电磁场，但计算需要选定时间并推进初值。本章把场强和电流按时间方向与空间切片分解，辨认哪些分量含时间导数、哪些限制初始数据。分量计数展示方程如何组织；约束能否随时间保持则留给下一章证明。
\end{chaptersummary}
\studyroute{这里略去 Lorentz Hodge 星的逐基计算和一般曲时空分叶。}{3+1 decomposition; Lorentzian Hodge star; spacetime foliation}{从 $\alpha\wedge\star\beta=\langle\alpha,\beta\rangle\vol_4$ 出发计算 $\star F$，逐项核对演化式中的符号。}
\section{为什么需要选定时间方向}
四维协变表达尤其适合理解几何不变量，但计算机上的时间推进或 PDE 初值理论需要指定一个初始空间超曲面 $\Sigma_0$ 和一族等时超曲面 $\Sigma_t$。在平直时空先取
\begin{equation}
M=\RR_t\times\Sigma,\qquad \Sigma=\RR^3.
\end{equation}
这一分解用于\emph{构造演化问题}；不同惯性观察者可分别选择自己的时间切片。

\section{时空外微分的拆分}
在当前乘积坐标中，$\partial_t$ 是时间坐标向量场，
$\dd t$ 是对偶坐标基中的一形式，满足 $\dd t(\partial_t)=1$。
对依赖时间的空间 $k$ 形式，在不随时间变化的空间坐标基下定义
\[
\alpha(t)=\sum_I\alpha_I(t,\bm x)\,\dd x^I,\qquad
\partial_t\alpha:=\sum_I(\partial_t\alpha_I)\,\dd x^I,
\]
其中 $I=(i_1<\cdots<i_k)$，$\dd x^I=\dd x^{i_1}\wedge\cdots\wedge\dd x^{i_k}$。
这里先对系数求时间导数，所得 $\partial_t\alpha$ 仍为空间 $k$ 形式。
将外微分的坐标表达按时间与空间分组，得到
\begin{equation}\label{eq:splitd}
\dd\alpha=\dd_\Sigma\alpha+\dd t\wedge(\partial_t\alpha).
\end{equation}
两项都是 $(k+1)$ 形式；第二项表示复合算子
$(\dd t\wedge\cdot)\circ\partial_t$ 的作用，而非一形式与向量场的外积。
例如 $\alpha=t\,\dd x$ 时，$\partial_t\alpha=\dd x$，
故 $\dd\alpha=\dd t\wedge\dd x$。
在当前固定乘积坐标中，$\partial_t\alpha$ 也等于沿 $\partial_t$ 的
Lie 导数 $\mathcal L_{\partial_t}\alpha$。
坐标基与外微分可参见
\href{https://www.davidtong.org/teaching/general-relativity/grhtml/S2}{Tong 讲义第 2.3.1、2.4.1 节}；
以 Lie 导数表述的分解见
\href{https://arxiv.org/abs/gr-qc/9701066}{Fecko 第 IV.B 节}。

\begin{askai}
从 $F=B-\dd t\wedge E$ 出发，利用式 \eqref{eq:splitd} 逐项展开 $\dd F$。
先对 $B,E$ 的系数求时间导数，再与 $\dd t$ 外积；请核对每一步的形式次数，
指出每个符号来自外积交换还是时间求导，并分别识别空间三形式与含 $\dd t$ 的部分。

进一步展开：把得到的两条形式方程翻译成三维向量方程，逐项核对符号。
\end{askai}
由反对称性和外微分的乘积法则，对任意空间形式 $\beta$：
\begin{equation}
\dd(\dd t\wedge\beta)=-\dd t\wedge\dd_\Sigma\beta.
\end{equation}
把 $F=B-\dd t\wedge E$ 代入：
\begin{align}
\dd F
 &=\dd_\Sigma B+\dd t\wedge\partial_tB
   +\dd t\wedge\dd_\Sigma E\\
 &=\underbrace{\dd_\Sigma B}_{\text{纯空间三形式}}
  +\dd t\wedge\underbrace{(\partial_tB+\dd_\Sigma E)}_{\text{空间二形式}}.
\end{align}
因此
\begin{equation}\label{eq:splitFaraday}
\boxed{\partial_tB+\dd_\Sigma E=0,\qquad \dd_\Sigma B=0.}
\end{equation}
前式是 Faraday 演化方程（3 个标量分量），后式是无磁单极约束（1 个标量分量）。

\section{Hodge 星如何恢复另外四条标量方程}
我们的度规与取向给出
\begin{equation}\label{eq:starF}
\star F=\starS E+\dd t\wedge\starS B.
\end{equation}
注意 $\starS B$ 是空间一形式。于是
\begin{align}
\dd\star F
 &=\dd_\Sigma\starS E
   +\dd t\wedge\bigl(\partial_t\starS E-\dd_\Sigma\starS B\bigr).
\end{align}
将 $J=\rho\vol_3-\dd t\wedge\starS j^\flat$ 代入 $\dd\star F=J$，得到
\begin{equation}\label{eq:splitAmpere}
\boxed{
\begin{aligned}
\partial_t\starS E-\dd_\Sigma\starS B&=-\starS j^\flat,\\
\dd_\Sigma\starS E&=\rho\vol_3.
\end{aligned}}
\end{equation}
将空间一形式与三维向量对应，就恢复了 \eqref{eq:maxwell3} 的 Amp\`ere--Maxwell 演化方程和 Gauss 电约束。

\section{为什么三形式恰好拆成 \texorpdfstring{$3+1$}{3+1} 个分量}
在每个选定的时空点，四维三形式可唯一写作
\begin{equation}\label{eq:3formsplit}
\Omega^3(M)\simeq
\bigl(\dd t\wedge\Omega^2(\Sigma)\bigr)\oplus\Omega^3(\Sigma).
\end{equation}
这是相对于时间切片的\emph{逐点向量空间分解}，各分量仍可依赖时间。计数上，
\begin{equation}
\binom43=\binom32+\binom33=3+1.
\end{equation}
于是每条四维三形式 Maxwell 方程，都分裂成 3 个演化分量和 1 个约束分量。两条合起来是\textbf{6 个演化分量加 2 个约束分量}。

\begin{center}
\begin{tikzpicture}[>=stealth,scale=1,every node/.style={font=\small}]
\node[draw,rounded corners,minimum width=33mm,minimum height=9mm] (a) at (0.75,0) {$\dd F=0$};
\node[draw,rounded corners,minimum width=33mm,minimum height=9mm] (b) at (8,0) {$\dd\star F=J$};
\node[draw,rounded corners,align=center,minimum width=30mm] (c) at (-1,-1.6) {3 个演化分量\\$\partial_t B+\dd_\Sigma E=0$};
\node[draw,rounded corners,align=center,minimum width=27mm] (d) at (2.5,-1.6) {1 个约束分量\\$\dd_\Sigma B=0$};
\node[draw,rounded corners,align=center,minimum width=36mm] (e) at (6,-1.6) {3 个演化分量\\Amp\`ere--Maxwell};
\node[draw,rounded corners,align=center,minimum width=27mm] (f) at (10,-1.6) {1 个约束分量\\Gauss 电定律};
\draw[->] (a)--(c);\draw[->] (a)--(d);\draw[->] (b)--(e);\draw[->] (b)--(f);
\end{tikzpicture}
\end{center}

\section{演化与约束的依赖关系}
\begin{enumerate}
\item \textbf{约束限定初值集合。}两个 Gauss 方程约束同一组电磁场变量。
\item \textbf{时间切片指定演化方向。}协变方程可相对于所选观察者写成初值问题。
\item \textbf{约束传播依赖微分恒等式。}其证明使用 $\dd^2=0$ 和源的连续性方程。
\end{enumerate}

\section{练习}
\begin{enumerate}
\item 按照 $F=B+E\wedge\dd t$ 逐步核对 \eqref{eq:splitFaraday} 中的符号。
\item 在二维空间加一维时间中，讨论 $\Omega^2$ 的时间/空间分量数目。
\item 解释 $\dd\star F=J$ 的约束中为什么出现 $\rho$，演化中为什么出现 $\bm j$。
\end{enumerate}

\chapter{约束传播、守恒恒等式与初值集合}
\label{ch:constraint}
\begin{chaptersummary}
上一章分离出演化方程和 Gauss 约束，本章追问：满足约束的初值在演化中是否仍然满足约束？答案来自旋度与散度的复形关系及电荷连续性。这个连续层面的相容性也为后面构造离散算子和检查散度误差提供标准。
\end{chaptersummary}
\studyroute{这里略去边界适定性与全局拓扑分析，集中证明光滑解的局部约束传播。}{constraint propagation; de Rham complex; discrete continuity equation}{给出满足和不满足 $D_hC_h=0$ 的矩阵例子，比较磁散度在一步更新后的变化。}
\section{约束误差的定义}
从此回到最便于工程计算的三维向量写法。只把旋度方程作为演化系统：
\begin{equation}\label{eq:evol}
\partial_t\bm B=-\curl\bm E,
\qquad
\partial_t\bm E=\curl\bm B-\bm j.
\end{equation}
定义
\begin{equation}
C_B=\Div\bm B,\qquad C_E=\Div\bm E-\rho.
\end{equation}
$C_B=C_E=0$ 才代表一个物理允许的初始场。

\section{经典层面的约束传播}
\begin{proposition}[光滑解的约束传播]\label{prop:propagate}
设 $\bm E,\bm B,\rho,\bm j$ 足够光滑，满足演化方程 \eqref{eq:evol}，并且 $\partial_t\rho+\Div\bm j=0$。则
\begin{equation}
\partial_tC_B=0,\qquad \partial_tC_E=0.
\end{equation}
特别地，若 $t=0$ 时 $C_B=C_E=0$，其后在光滑解存在的时间区间内仍然成立。
\end{proposition}
\begin{proof}
对磁场演化方程取散度，由 $\Div\curl=0$ 得
\begin{equation}
\partial_t C_B=-\Div\curl\bm E=0.
\end{equation}
对电场演化方程取散度，使用连续性方程，有
\begin{align}
\partial_tC_E
&=\Div(\curl\bm B-\bm j)-\partial_t\rho\\
&=-\Div\bm j-\partial_t\rho=0.
\end{align}
\end{proof}

\section{几何层面的解释}
\begin{studentinput}
假设初值满足 $\Div\bm B_0=0$。请从磁场演化方程推导该约束的时间导数，并解释证明中哪一步对应 $\dd^2=0$。电 Gauss 约束还需要源项满足什么条件？

\textbf{你的答案：}\vspace{2cm}
\end{studentinput}
命题 \ref{prop:propagate} 的结构基础是 de Rham 复形
\begin{equation}\label{eq:derham}
\Omega^0\xrightarrow{\dd}\Omega^1\xrightarrow{\dd}\Omega^2\xrightarrow{\dd}\Omega^3,
\qquad \dd\circ\dd=0.
\end{equation}
对 $\partial_tB+\dd_\Sigma E=0$ 再取 $\dd_\Sigma$，得
\begin{equation}
\partial_t(\dd_\Sigma B)=-\dd_\Sigma^2E=0.
\end{equation}
电约束同理，但还需要 $\dd J=0$。\textbf{约束传播是四维几何结构的时间投影。}

\section{初值与边界的相容条件}
给定 $\bm E(0)=\bm E_0,\bm B(0)=\bm B_0$ 时，初值须满足
\begin{equation}\label{eq:initConstraints}
\Div\bm B_0=0,
\qquad \Div\bm E_0=\rho(0,\cdot).
\end{equation}
在全空间上，还需适宜的无穷远衰减或有限能量假设；在有界区域上，需附加与演化兼容的边界条件。比如封闭区域总电荷与电通量之间必须满足积分 Gauss 定律。边界条件还需与能量通量及入射特征相容。

\section{离散复形与约束传播}
\begin{engineernote}
网格的点、边、面、体及其取向给出关联矩阵 $G,C,D$。在相容的关联关系中，$CG=0$、$DC=0$ 是拓扑恒等式；几何和材料信息则进入质量矩阵或离散 Hodge 型算子。实际代码还需检查场变量所在的网格及边界处理。
\end{engineernote}

\begin{askai}
请为一个取向明确的四面体编号顶点、边和面，写出顶点--边关联矩阵 $G$、边--面关联矩阵 $C$ 和面--体关联矩阵 $D$，逐项验证 $CG=0$ 和 $DC=0$，解释每个符号。

进一步展开：反转一条边的取向后，哪些矩阵行或列改变？恒等式为什么仍成立？
\end{askai}
真实数值计算的导数由离散矩阵代替。若离散散度 $D_h$ 与离散旋度 $C_h$ 未满足
\begin{equation}\label{eq:discretecomplex}
D_hC_h=0,
\end{equation}
则离散演化可能产生非零磁散度。电约束还需要电流/电荷更新满足\emph{相同离散算子}下的连续性方程。第 6 章和第 7 章将据此讨论保约束离散。

\section{练习}
\begin{enumerate}
\item 取 $\rho=0,\bm j=0$，给出一个 $\Div\bm B_0\neq0$ 的光滑初值。证明仅靠演化方程，它的磁约束错误不会自行消失。
\item 对有边界的体积 $V$，把 $\Div\bm E=\rho$ 积分，说明初值和边界电通量须满足什么关系。
\item 从离散恒等式 $D_hC_h=0$ 推出最简单的离散磁约束传播公式。
\end{enumerate}

\chapter{对称双曲系统：特征、能量与适定性}
\label{ch:hyperbolic}
\begin{chaptersummary}
约束确定了允许的初始场；本章转向这些场如何传播。将 Maxwell 演化方程写成对称双曲系统后，特征值给出传播速度，能量估计控制解对数据的依赖。它们为守恒律、数值通量和时间步长分析提供依据，同时须结合边界与存在性讨论适定性。
\end{chaptersummary}
\studyroute{这里略去存在性定理的证明以及入射特征边界条件的完整分类。}{symmetric hyperbolic systems; energy estimate; maximally dissipative boundary conditions}{对 $A(\bm n)$ 分解入射、出射与零速子空间，说明边界数据如何进入能量通量。}
\section{把 Maxwell 演化式写成一阶矩阵系统}
令
\begin{equation}
U=\mat{\bm E\\\bm B}\in\RR^6,
\qquad S=\mat{-\bm j\\\bm0}.
\end{equation}
定义 $3\times3$ 的反对称矩阵 $C(\bm n)$，使
\begin{equation}
C(\bm n)\bm v=\bm n\times\bm v,
\qquad C(\bm n)^\mathsf T=-C(\bm n).
\end{equation}
令
\begin{equation}\label{eq:An}
A(\bm n)=
\begin{pmatrix}
0&-C(\bm n)\\
C(\bm n)&0
\end{pmatrix}.
\end{equation}
于是 $A(\bm n)^\mathsf T=A(\bm n)$，并且取 $A^i=A(\bm e_i)$ 后，
\begin{equation}\label{eq:matrixhyper}
\boxed{\partial_tU+\sum_{i=1}^3A^i\partial_iU=S.}
\end{equation}
一个一阶系统若能在正定内积下把主部矩阵同时对称化，就称\emph{对称双曲系统}。此例已经是对称的，对称化矩阵可直接取 $I_6$。

\section{沿任意方向的传播速度}
\begin{askai}
取传播方向 $\bm n=(1,0,0)$，从本章的 Maxwell 主部矩阵 $A(\bm n)$ 出发求特征值与特征向量，区分横向传播分量和零速纵向分量；再说明旋转坐标如何推广到任意单位方向。

进一步展开：在真空、周期边界下，对该系统逐步推导能量估计，并标明边界通量在哪里消失。
\end{askai}
取单位向量 $\bm n$，其法向主部由 $A(\bm n)$ 决定。将 $\bm E,\bm B$ 拆成沿 $\bm n$ 的纵向分量与垂直 $\bm n$ 的横向分量。利用
\begin{equation}
C(\bm n)^2\bm v=\bm n(\bm n\cdot\bm v)-\bm v,
\end{equation}
可知 $A(\bm n)$ 的特征值是
\begin{equation}\label{eq:eigenvalues}
\boxed{-1,-1,0,0,1,1.}
\end{equation}
$\pm1$ 是横向电磁波的传播速度，重复次数各为 2；两个零特征值对应纵向分量。纵向分量的取值受 Gauss 约束限制。由于真空单位选取 $c=1$，实际量纲下传播速度为 $\pm c$。

\section{电磁能量估计}
将 $\partial_t\bm E=\curl\bm B-\bm j$ 与 $\partial_t\bm B=-\curl\bm E$ 分别点乘 $\bm E,\bm B$，利用向量恒等式
\begin{equation}
\Div(\bm E\times\bm B)=\bm B\cdot\curl\bm E-\bm E\cdot\curl\bm B,
\end{equation}
得局部能量平衡
\begin{equation}\label{eq:poynting}
\boxed{\partial_t\left(\tfrac12|\bm E|^2+\tfrac12|\bm B|^2\right)
+\Div(\bm E\times\bm B)=-\bm j\cdot\bm E.}
\end{equation}
对空间区域 $\Omega$ 积分，
\begin{equation}\label{eq:energyintegral}
\frac{\dd}{\dd t}\frac12\int_\Omega(|\bm E|^2+|\bm B|^2)\,\dd x
=-\int_{\partial\Omega}(\bm E\times\bm B)\cdot\bm n\,\dd S
-\int_\Omega\bm j\cdot\bm E\,\dd x.
\end{equation}
全空间快速衰减或周期边界下，边界通量消失。当 $\bm j=0$ 时，$L^2$ 能量守恒；有外源时，利用 Cauchy--Schwarz 不等式可以估计能量增长。

\begin{proposition}[全空间光滑解的稳定性估计]
设 $U$ 为 \eqref{eq:matrixhyper} 的光滑、充分衰减解，则
\begin{equation}\label{eq:energyestimate}
\norm{U(t)}_{L^2}\leq\norm{U(0)}_{L^2}+\int_0^t\norm{S(s)}_{L^2}\,\dd s.
\end{equation}
\end{proposition}
\begin{proof}
由对称矩阵主部和分部积分，$\tfrac12\frac{\dd}{\dd t}\norm U_{L^2}^2=\ip{U}{S}_{L^2}$。对 $\norm U_{L^2}>0$ 的时间段除以范数，积分得到结论；范数可能为零的情形可用正则化 $\sqrt{\norm U^2+\varepsilon^2}$ 处理。
\end{proof}
\begin{remark}
能量估计为稳定性与唯一性提供控制。存在性可通过常系数 Fourier 方法、半群方法或对称双曲系统理论建立；有界区域还需选择适当的边界条件并控制能量通量。
\end{remark}

\section{一维横向特征变量}
考虑只依赖 $x,t$ 的真空场，取横向分量 $E_y,B_z$。方程退化为
\begin{equation}\label{eq:1dwave}
\partial_tE_y+\partial_xB_z=0,
\qquad \partial_tB_z+\partial_xE_y=0.
\end{equation}
令 $w_+=E_y+B_z,\ w_-=E_y-B_z$，便有
\begin{equation}
\partial_tw_++\partial_xw_+=0,
\qquad
\partial_tw_- -\partial_xw_-=0.
\end{equation}
因此
\begin{equation}
w_+(x,t)=w_+(x-t,0),\qquad
w_-(x,t)=w_-(x+t,0).
\end{equation}
该分解展示双曲系统中信号沿有限速度特征线传播的性质。

\section{练习}
\begin{enumerate}
\item 验证 $A(\bm n)$ 为对称矩阵。
\item 对 $\bm n=\bm e_x$，直接计算 $A(\bm n)$ 的 6 个特征值。
\item 若初始 $E_y=B_z=f(x)$，其后沿哪个方向传播？若 $E_y=-B_z=f(x)$ 呢？
\item 说明为什么无穷远能量条件与有边界区域的入射特征边界条件在逻辑上不同。
\end{enumerate}

\chapter{从双曲系统到守恒律与弱解}
\label{ch:conservation}
\begin{chaptersummary}
上一章的双曲结构揭示了波的传播；本章将演化方程写成通量形式，并用积分意义处理可能不光滑的解。控制体平衡与跳跃条件把连续方程连接到界面信息交换，为有限体积和 DG 数值通量做准备。守恒型与能量平衡是相关但不同的性质。
\end{chaptersummary}
\studyroute{这里略去分布理论及一般非线性熵解理论，保留测试函数、通量平衡和跳跃关系。}{weak solution; Rankine--Hugoniot condition; entropy pair}{将分段常数行波代入分布意义的守恒律，推导 $s[U]=A(\bm n)[U]$。}
\section{守恒型表达}
由于 $A^i$ 为常矩阵，定义第 $i$ 个空间方向的通量
\begin{equation}\label{eq:flux}
\mathcal F^i(U)=A^iU
=\mat{-\bm e_i\times\bm B\\ \bm e_i\times\bm E}.
\end{equation}
则
\begin{equation}\label{eq:conservation}
\boxed{\partial_tU+\sum_{i=1}^3\partial_i\mathcal F^i(U)=S.}
\end{equation}
这里的“守恒型”指演化方程具有散度通量结构；有源项时，各场分量满足含源的平衡关系，Gauss 约束仍须单独考虑。

\section{控制体积分与有限体积法的起点}
对固定控制体 $K$ 积分，使用 Gauss 定理得到
\begin{equation}\label{eq:finitevolumeintegral}
\frac{\dd}{\dd t}\int_K U\,\dd x
+\int_{\partial K}\mathcal F(U)\cdot\bm n\,\dd S
=\int_KS\,\dd x.
\end{equation}
式中 $\mathcal F(U)\cdot\bm n:=\sum_i n_i\mathcal F^i(U)$ 为\emph{六维状态向量通量}。这就是有限体积更新公式的连续原型：网格单元内部平均值的变化，由所有外法向面上的通量和体积源项决定。

\section{弱形式与跳跃条件}
\begin{studentinput}
对于一维守恒律 $\partial_tU+\partial_x\mathcal F(U)=0$，假设解在 $x=st$ 两侧为常值。请用跨越界面的控制体说明为什么跳跃需要满足 $s[U]=[\mathcal F]$，并指出这一条件如何约束界面传播速度。

\textbf{你的答案：}\vspace{2cm}
\end{studentinput}
对低正则性解，使用积分意义下的方程。令 $\varphi\in C_c^\infty((0,T)\times\Omega;\RR^6)$，把 \eqref{eq:conservation} 乘以测试函数并积分分部：
\begin{equation}\label{eq:weakinternal}
\int_0^T\!\int_\Omega
\left(U\cdot\partial_t\varphi+
\sum_i\mathcal F^i(U)\cdot\partial_i\varphi+S\cdot\varphi\right)\dd x\,\dd t=0.
\end{equation}
若测试函数允许接触 $t=0$ 且在 $t=T$ 为零，则初值项出现为
\begin{equation}\label{eq:weakinit}
\int_0^T\!\int_\Omega
\left(U\cdot\partial_t\varphi+
\sum_i\mathcal F^i(U)\cdot\partial_i\varphi+S\cdot\varphi\right)\dd x\,\dd t
+\int_\Omega U_0\cdot\varphi(0,x)\,\dd x=0,
\end{equation}
这里假定测试函数的空间支集远离边界；两式采用整体积分分部后乘以 $-1$ 的符号约定。

在无奇异源的移动跳跃面上，若面法向为 $\bm n$、传播速度为 $s$，Rankine--Hugoniot 条件为
\begin{equation}\label{eq:RH}
s[U]=[\mathcal F(U)\cdot\bm n]=A(\bm n)[U],
\end{equation}
其中 $[U]=U_R-U_L$。因而非零跳跃只能沿对应特征速度传播。Maxwell 的线性系统没有 Euler 方程那样的非线性激波生成机制，但跃变传播和波分解仍然是数值通量的基础。

\section{守恒型与能量平衡}
$U$ 的六个分量满足矢量守恒型表达，而 $\tfrac12|U|^2$ 满足另一条\emph{能量平衡律} \eqref{eq:poynting}。后者在数值分析中可视为一个凸熵（entropy）及其通量：
\begin{equation}
\eta(U)=\tfrac12 U^\mathsf TU,\qquad
q^i(U)=\tfrac12 U^\mathsf TA^iU=(\bm E\times\bm B)_i.
\end{equation}
这种线性对称系统的“熵”是 PDE 术语，不等同于热力学熵；它用于稳定性分析与选择合适数值通量。

\section{练习}
\begin{enumerate}
\item 验证 $\mathcal F^x(U)$ 的前 3 个分量与 $-\bm e_x\times\bm B$ 一致。
\item 对一维系统 \eqref{eq:1dwave}，推导 Rankine--Hugoniot 条件并检查速度 $\pm1$。
\item 用对称矩阵性质验证 $q^i(U)$ 的公式。
\end{enumerate}


% 5.5 是在保留原第 6 章编号前提下插入的独立讲授单元。
\begingroup
\chapter*{第 5.5 章：从微分方程到算法——空间离散、数值通量与时间推进}
\phantomsection
\addcontentsline{toc}{chapter}{第 5.5 章：从微分方程到算法——空间离散、数值通量与时间推进}
\markboth{第 5.5 章 数值求解的三个选择}{第 5.5 章 数值求解的三个选择}
\setcounter{section}{0}
\renewcommand{\thesection}{5.5.\arabic{section}}
\renewcommand{\thesubsection}{5.5.\arabic{section}.\arabic{subsection}}
\renewcommand{\theHsection}{bridge.\arabic{section}}
\setcounter{equation}{0}
\renewcommand{\theequation}{5.5.\arabic{equation}}
\renewcommand{\theHequation}{bridge.\arabic{equation}}
\label{ch:bridge}

\begin{chaptersummary}
守恒律和弱形式仍是连续问题。本章明确通向算法的三个选择：空间自由度、界面数值通量和时间推进，并说明稳定性与长期结构取决于它们的配合。这一框架连接后续的有限体积、散度保持与高阶 DG；单独选择某一个方法并不能决定完整算法的全部性质。
\end{chaptersummary}
\studyroute{这里略去高阶方法的全部系数表与阶条件证明；新增 ADER、CERK、LTS 和 IMEX 小节给出计算结构与最小算例。}{method of lines; ADER predictor; continuous RK; conservative LTS; additive Runge--Kutta}{区分空间残差、步内预测、时间调度和刚性分拆，列出这些步骤各自需要的数据。}

\noindent\textbf{这一章的任务。}
从守恒型和弱形式构造计算程序，需要明确三个选择：
\textbf{空间中用什么有限个自由度表示场？单元边界如何交流信息？如何从 $t_n$ 推进到 $t_{n+1}$？}
三者相互关联。本章介绍共同框架，第 6 章讨论有限体积，
第 7 章介绍从 Sobolev 弱形式出发的 Galerkin--DG 路线。

\section{线法与空间半离散}
以 $\partial_t U+\nabla\cdot\mathcal F(U)=S$ 为例，
空间离散（有限差分、有限体积、有限元、DG、离散微分形式等）把连续未知场 $U(t,x)$
变成有限维系数 $u_h(t)\in\RR^N$，通常得到
\begin{equation}\label{eq:mol}
M_h\dot u_h=R_h(u_h,t),\qquad
\dot u_h=L_h(u_h,t):=M_h^{-1}R_h(u_h,t).
\end{equation}
$M_h$ 是质量矩阵，$R_h$ 含体积分、界面通量、边界条件与源项。这叫\emph{线法}
（method of lines）。另一些方法同时离散时间、空间；工程中线法是理解算法结构的好起点。

\begin{center}
\begin{tikzpicture}[>=stealth,node distance=6mm,
 every node/.style={font=\small},
 stage/.style={draw,rounded corners,align=center,text width=3.2cm,minimum height=13mm}]
\node[stage] (pde) {PDE 与弱形式\\$\partial_t U+\nabla\cdot\mathcal F(U)=S$};
\node[stage,right=of pde] (space) {空间离散与界面通量\\$M_h\dot u_h=R_h(u_h)$};
\node[stage,right=of space] (time) {时间推进\\$u_h^{n+1}=\Psi_{\Delta t}(u_h^n)$};
\draw[->,thick] (pde)--(space);
\draw[->,thick] (space)--(time);
\end{tikzpicture}
\end{center}

\paragraph{什么是数值通量？}
相邻单元可使用不同的近似，界面因而存在两个迹 $U^-$、$U^+$。
它们必须通过\emph{同一个}数值通量连接。
记该通量为 $\widehat{\mathcal F}(U^-,U^+;n)$，
要求它至少满足一致性与相邻单元法线反向的守恒性：
\begin{equation}
\widehat{\mathcal F}(U,U;n)=\mathcal F(U)\cdot n,
\end{equation}
以及
\begin{equation}
\widehat{\mathcal F}(U^-,U^+;n)
=-\widehat{\mathcal F}(U^+,U^-;-n).
\end{equation}
中心通量将两边信息对称平均；上风通量依据传播方向附加适当耗散。
\textbf{通量选择属于空间算子设计；时间积分器负责推进所得的半离散系统。}

\section{常见时间推进：精度、稳定性、成本}
\begin{checkpoint}
对于半离散方程 $\dot u=L_hu$，为什么空间算子守能仍不足以保证显式 Euler 稳定？请用纯虚特征值检验一步放大因子，并说明时间步长与稳定域的关系。
\end{checkpoint}
对 $\dot u=L_h(u)$，显式 Euler 为
$u^{n+1}=u^n+\Delta t\,L_h(u^n)$；
改用经典四阶 Runge--Kutta（RK4）可获得四阶时间截断精度，但需四次右端计算。
对于简单的半离散标量模态 $\dot y=\lambda_h y$，每种时间方法都对应放大因子
$y^{n+1}=R(z)y^n$，$z=\Delta t\lambda_h$：
\begin{align}
R_{\mathrm{FE}}(z)&=1+z,\\
R_{\mathrm{RK4}}(z)&=1+z+\frac{z^2}{2}+\frac{z^3}{6}+\frac{z^4}{24}.
\end{align}
必须在离散算子谱与时间方法稳定域之间检查相容性。
例如，纯虚谱上的显式 Euler 满足 $|1+\mathrm i\alpha|>1$（$\alpha\ne0$），
因此显式 Euler 在这一固定步长下放大纯振荡模态。

\paragraph{CFL 条件及其依赖因素。}
显式双曲求解器一般需要
\begin{equation}\label{eq:cflgeneral}
\Delta t\lesssim \frac{h}{c\,C(p,\text{网格},\text{通量},\text{时间格式})},
\end{equation}
其中 $c$ 是特征速度，$p$ 是多项式次数。高阶 DG 的局部微分算子谱会随 $p$
扩张，故提高空间阶数通常使显式稳定步长更严格。
隐式方法可放宽线性稳定性限制；步长选择仍需同时考虑精度、求解成本及约束相容性。

\section{保辛与长时间电磁演化}\label{sec:bridge-symp}
\begin{askai}
对简谐振子分别写出显式 Euler 与隐式中点法的一步矩阵，计算它们对辛形式及二次能量的影响，再比较数值频率与真实频率。

进一步展开：对一般非线性 Hamilton 系统，为什么保辛不能单独推出逐步精确守能？
\end{askai}
在无源、无耗、适当边界条件下，Maxwell 系统有 Hamilton/Poisson 结构。
若一种\emph{兼容空间离散}产生有限维系统
\begin{equation}\label{eq:discrete-hamiltonian}
\dot u_h=P_h\nabla H_h(u_h),\qquad
P_h^\mathsf T=-P_h,\quad
H_h(u_h)=\tfrac12u_h^\mathsf TK_hu_h,
\quad K_h=K_h^\mathsf T>0,
\end{equation}
则因反对称性 $\dot H_h=0$。在 Maxwell 的非正则 Poisson 写法中，
$P_h$ 可能\emph{退化}，散度约束表现为 Casimir 或不变量；
此时应在辛叶或适当约束变量上讨论正则坐标。
若限制在合适的非退化辛叶上，或通过势变量和约束处理获得正则 Hamilton 系统，
可使用标准保辛方法。

\paragraph{隐式中点法。}
对 $\dot u_h=A_hu_h$，取
\begin{equation}\label{eq:midpoint}
\frac{u_h^{n+1}-u_h^n}{\Delta t}
=A_h\frac{u_h^{n+1}+u_h^n}{2}.
\end{equation}
若存在正定能量矩阵 $Q_h$，使
$A_h^\mathsf TQ_h+Q_hA_h=0$，则将
\eqref{eq:midpoint} 左乘
$(u_h^{n+1}+u_h^n)^\mathsf TQ_h/2$，立即得到
\begin{equation}\label{eq:midpoint-energy}
\frac12(u_h^{n+1})^\mathsf TQ_hu_h^{n+1}
=\frac12(u_h^n)^\mathsf TQ_hu_h^n.
\end{equation}
这是\textbf{线性、无耗、离散能量内积兼容}情况下的精确二次能量保持；
中点法在正则 Hamilton 系统上也保辛。
式 \eqref{eq:midpoint-energy} 使用线性斜伴随条件；非线性、材料耗散、PML 或外部做功需要相应的离散能量平衡分析。

\paragraph{St\"ormer--Verlet / leapfrog 与 Yee--FDTD。}
振子可采用交错推进
\begin{align}\label{eq:leapfrogosc}
p^{n+1/2}&=p^{n-1/2}-\Delta t\,\omega^2q^n,\\
q^{n+1}&=q^n+\Delta t\,p^{n+1/2}.
\end{align}
在正则的振子模型中，这类分裂方法具有保辛结构。
电磁计算中常见的 Yee 交错时间更新，使用相邻半时间层的 $E$、$B$
（具体放置取决于约定），其\emph{无源、常系数、兼容离散且边界处理适当}的
线性更新与交错 Hamilton 推进密切相关。
Yee 网格的正则保辛性取决于介质、边界和具体修正；第 6 章重点使用其离散散度恒等式。

\paragraph{保辛与数值色散}
对纯振荡 $\dot y=\mathrm i\omega y$，中点法的放大因子为
\begin{equation}
R_{\rm MP}=\frac{1+\mathrm i\omega\Delta t/2}{1-\mathrm i\omega\Delta t/2},
\qquad |R_{\rm MP}|=1,
\end{equation}
但它给出的数值角频率是
\begin{equation}\label{eq:midpointdisp}
\omega_{\Delta t}=\frac2{\Delta t}
\arctan\!\left(\frac{\omega\Delta t}{2}\right)
\ne\omega\quad(\Delta t>0,\;\omega\ne0).
\end{equation}
因此\emph{振幅保持}与\emph{相位精度}需要分别评估。
长时间电磁波计算既要关心几何/能量稳定，也要控制时空离散的总色散误差；
第 7 章将专门检验高阶 DG 的\emph{空间}色散。

\section{ADER：局部时空预测与校正}\label{sec:ader}
式 \eqref{eq:mol} 给出空间残差，而高阶更新还需要描述一步之内的演化。
ADER（Arbitrary high-order DERivatives）通过时间展开或局部时空预测构造这类信息。
对常系数平流 $\partial_tu+a\partial_xu=0$，经典 Taylor 构造使用
\begin{equation}\label{eq:ader-taylor}
\partial_t^m u=(-a)^m\partial_x^m u,\qquad
q_K(x,t_n+\tau)=\sum_{m=0}^{r}
\frac{\tau^m}{m!}(-a)^m\partial_x^m u_h(x,t_n).
\end{equation}
这是 Cauchy--Kovalevskaya 过程的线性例子。对一般系统，现代 ADER 也可由单元内的时空弱问题求得
$q_K(\bm x,t)$，随后把预测值代入数值通量，完成与邻居耦合的校正。
以无源有限体积更新为例，
\begin{equation}\label{eq:ader-corrector}
|K|(\bar U_K^{n+1}-\bar U_K^n)
=-\sum_{f\subset\partial K}\int_{t_n}^{t_{n+1}}\int_f
\widehat{\mathcal F}(q_K,q_{K_f};\bm n_f)\,\dd S\,\dd t.
\end{equation}
DG 更新还包含基函数加权的体积项。局部预测器的可计算性来自单元内信息；
守恒耦合来自共享界面通量。预测阶数、积分阶数与空间精度共同决定全离散阶数，
显式双曲校正仍需相应 CFL 条件。局部预测与时间自适应构造参见
\href{https://doi.org/10.1016/j.cma.2014.07.019}{Dumbser (2014)}。

\section{CERK：步内连续时间近似}\label{sec:cerk}
连续显式 Runge--Kutta（continuous explicit Runge--Kutta, CERK）用阶段导数构造整个时间步上的近似。
对自治系统 $\dot u=L_h(u)$，写成
\begin{equation}\label{eq:cerk}
\begin{aligned}
k_i&=L_h\left(u^n+\Delta t\sum_{j<i}a_{ij}k_j\right),\\
u_h(t_n+\theta\Delta t)&\approx u^n+
\Delta t\sum_{i=1}^{s}b_i(\theta)k_i,\qquad 0\leq\theta\leq1.
\end{aligned}
\end{equation}
$b_i(0)=0$，$b_i(1)=b_i$；非自治问题还需在阶段时刻评价右端。
连续扩展具有自己的阶条件。以二阶 Heun 法为例，
\begin{equation}\label{eq:cerk-heun}
k_1=L_h(u^n),\quad k_2=L_h(u^n+\Delta t k_1),\qquad
b_1(\theta)=\theta-\tfrac12\theta^2,\quad
b_2(\theta)=\tfrac12\theta^2.
\end{equation}
在光滑条件下，该多项式给出二阶步内近似；取 $\theta=1$ 恢复 Heun 更新。
高阶 CERK 的系数应由连续阶条件确定，参见
\href{https://doi.org/10.1137/0913084}{Owren--Zennaro (1992)}。
连续扩展可用于输出与时间耦合；将它用于单元局部预测时，阶段计算也须按局部问题组织，
再由校正通量加入邻居影响，见
\href{https://congress2.cimne.com/eccomas/proceedings/cfd2010/papers/01008.pdf}{Altmann--Gassner--Munz (2010)}。

\section{LTS：时间分层与共享通量}\label{sec:lts}
网格尺寸、材料波速和多项式次数造成不同单元的稳定步长差异。
局部时间步进（local time stepping, LTS）据此分配 $\Delta t_K$，并在共同时间层同步。
例如粗单元采用 $\Delta T$，细单元采用 $\Delta T/2$；在固定法向下，
共享面上的积分必须满足
\begin{equation}\label{eq:lts-flux}
\Phi_f[t_n,t_n+\Delta T]
=\Phi_f[t_n,t_n+\Delta T/2]
+\Phi_f[t_n+\Delta T/2,t_n+\Delta T],
\end{equation}
其中 $\Phi_f[a,b]=\int_a^b\int_f\widehat{\mathcal F}\,\dd S\,\dd t$。
实现时累积同一组子区间通量，并向相邻单元施加相反符号。
ADER 或 CERK 可提供各积分时刻的界面预测值；时间分层、数据可用性、积分阶数与同步规则
仍是 LTS 的组成部分。真实时间 LTS 需同时保持时间精度和界面守恒。
对于 Maxwell，电荷更新与磁场边环量也需时间相容，才能维持相应 Gauss 约束。

\begin{askai}
\textbf{略去内容与关键词：}continuous RK order conditions、local predictor、conservative LTS；这里略去高阶系数表与调度实现。

对一维平流的两个相邻单元采用 $2:1$ 时间步长比，使用式 \eqref{eq:cerk-heun} 提供步内值，
列出通量积分、细步更新与粗步同步的顺序。证明共享通量在两单元总量中抵消。

进一步展开：将粗侧状态冻结为 $u^n$，通过步长减半实验比较时间阶数。
\end{askai}

\section{IMEX：传播与刚性源项的分拆}\label{sec:imex}
ADER、CERK 与 LTS 处理时间近似及局部步长差异。
当材料响应或碰撞时间远短于目标演化时间时，还会出现刚性。
隐显式方法（implicit--explicit, IMEX）将右端拆为
\begin{equation}\label{eq:imex-split}
\dot u_h=F_h(u_h)+G_h(u_h),
\end{equation}
其中 $F_h$ 显式处理，$G_h$ 隐式处理。分拆依据是谱尺度与求解成本，
例如对旋度传播显式推进，对快速局部耗散隐式推进。
一个自治加性 RK 格式可写成
\begin{equation}\label{eq:imex-rk}
\begin{aligned}
Y_i&=u^n+\Delta t\sum_{j<i}\widetilde a_{ij}F_h(Y_j)
           +\Delta t\sum_{j\le i}a_{ij}G_h(Y_j),\\
u^{n+1}&=u^n+\Delta t\sum_i\widetilde b_iF_h(Y_i)
              +\Delta t\sum_i b_iG_h(Y_i).
\end{aligned}
\end{equation}
$\widetilde A$ 严格下三角，$A$ 可取对角隐式形式。
两个系数表须满足加性阶条件，见
\href{https://doi.org/10.1016/S0168-9274(97)00056-1}{Ascher--Ruuth--Spiteri (1997)}。
若 $G_h$ 仅含局部材料变量且质量矩阵按单元分块，阶段求解可限制在各块内；
若隐式部分包含空间微分耦合，则需要相应的全局线性或非线性求解。

\paragraph{导电 Maxwell 示例。}
本段恢复正的常数 $\varepsilon,\mu$，引入欧姆电流 $\bm j=\sigma\bm E$，$\sigma\geq0$。
于是
\begin{equation}\label{eq:conductive-maxwell}
\varepsilon\partial_t\bm E=\curl\bm H-\sigma\bm E,\qquad
\mu\partial_t\bm H=-\curl\bm E.
\end{equation}
电场松弛时间为 $\tau_\sigma=\varepsilon/\sigma$（$\sigma>0$）。
纯耗散模态 $\dot E=-E/\tau_\sigma$ 的显式 Euler 稳定步长为 $\Delta t\leq2\tau_\sigma$。
一阶 IMEX Euler 展示最基本的分拆：
\begin{equation}\label{eq:imex-euler-maxwell}
\begin{aligned}
(\varepsilon+\sigma\Delta t)\bm E^{n+1}
 &=\varepsilon\bm E^n+\Delta t\,\curl_h\bm H^n,\\
\mu\bm H^{n+1}&=\mu\bm H^n-\Delta t\,\curl_h\bm E^n.
\end{aligned}
\end{equation}
这里假定系数按点或块相乘；一般有限元形式需使用材料质量矩阵。
纯耗散模态的放大因子变为 $(1+\Delta t/\tau_\sigma)^{-1}$。
传播部分仍受显式稳定域限制：当 $\sigma=0$ 且空间算子具有纯虚谱时，
这一一阶例子退化为不适于无耗波的显式 Euler。
实际波动求解应选择虚轴稳定性合适的显式表，并检查完整 IMEX 放大矩阵。

在无边界能量通量时，连续模型满足
\begin{equation}\label{eq:conductive-energy}
\frac{\dd}{\dd t}\frac12\int_\Omega
(\varepsilon|\bm E|^2+\mu|\bm H|^2)\,\dd x
=-\int_\Omega\sigma|\bm E|^2\,\dd x.
\end{equation}
因此离散验证目标是含物理耗散的能量平衡。
电 Gauss 约束还要求电荷更新使用与电场方程相同时间权重的电流；
对于式 \eqref{eq:imex-euler-maxwell}，应相容地使用 $\sigma\bm E^{n+1}$。
若研究未解析松弛极限，需进一步检验极限模型、渐近保持性与刚性条件下的阶数。

\begin{askai}
\textbf{略去内容与关键词：}IMEX Runge--Kutta、additive order conditions、stiff accuracy、asymptotic preserving；这里略去高阶系数构造与刚性极限证明。

对式 \eqref{eq:conductive-maxwell} 的一维横向 Fourier 模态构造 $2\times2$ 放大矩阵，
比较显式 Euler 与式 \eqref{eq:imex-euler-maxwell} 在不同 $\sigma\Delta t/\varepsilon$ 下的谱半径、
能量和相位误差。区分消除耗散步长限制与稳定传播这两个条件。

进一步展开：选取文献中一组明确的 IMEX 系数，报告出处，并以时间收敛实验检查其实际阶数。
\end{askai}

\section{求解器设计的检查表}
\begin{enumerate}
\item 先规定要保持的量：电荷连续性、两类 Gauss 约束、离散能量，还是长时间相位精度？
\item 再选空间自由度与通量：单元平均、边/面自由度、连续或不连续的局部多项式。
\item 随后选时间格式：稳定域、计算成本、时间阶数，以及是否需要特定保结构性质。
\item 最后分别验证：空间收敛、时间收敛、数值散度、能量平衡、相位/振幅误差。
\item 对 LTS 记录各时间层更新次数与界面通量平衡；对 IMEX 分别检查传播 CFL、刚性衰减与阶段求解成本。
\end{enumerate}
\section{练习}
\begin{enumerate}
\item 写出显式 Euler 用于 $\dot y=\mathrm i\omega y$ 的增长因子，解释为何不适合纯振荡模式。
\item 用式 \eqref{eq:midpoint-energy} 的条件证明隐式中点法的离散二次能量恒等式。
\item 由式 \eqref{eq:midpointdisp} 展开 $\omega_{\Delta t}$ 的前三个非零 Taylor 项，说明时间积分本身也引入色散。
\end{enumerate}
\endgroup

\chapter{数值求解：双曲通量与散度约束}
\label{ch:numerics}
\begin{chaptersummary}
前一章区分了空间通量与时间推进，本章把它们用于 Maxwell 型双曲方程，并检查连续散度约束在网格上如何演化。有限体积通量负责相邻单元的信息交换，Yee 型交错结构则利用离散旋度与散度的相容性保持磁约束。两种要求需要共同纳入求解器设计。
\end{chaptersummary}
\studyroute{这里略去完整三维 FDTD 的材料、边界和电流沉积实现。}{Yee scheme; constrained transport; charge-conserving current deposition; GLM}{构造一个立方体的边--面--体关联矩阵，说明共享边电场如何保证磁通总和不变。}
\section{半离散有限体积格式}
将区域划分为控制体 $K$，定义单元平均
\begin{equation}
\bar U_K(t)=\frac1{|K|}\int_K U(t,x)\,\dd x.
\end{equation}
对每一单元面 $f\subset\partial K$ 选择一致的数值法向通量 $\widehat{\mathcal F}(U^-,U^+;\bm n_f)$，得到
\begin{equation}\label{eq:semidiscrete}
\frac{\dd\bar U_K}{\dd t}=
-\frac1{|K|}\sum_{f\subset\partial K}|f|\,
\widehat{\mathcal F}(U_K^-,U_{K_f}^+;\bm n_f)+\bar S_K.
\end{equation}
相邻单元共用同一数值面通量（取相反法向后变号），因而内部通量在全局求和中相互抵消，表现出有限体积法的局部保守性。

\section{两个可实现的数值通量}
真空 Maxwell 的法向 Jacobian 为 $A(\bm n)$，最大特征速度绝对值为 1。最容易实现的是 Rusanov（local Lax--Friedrichs）通量：
\begin{equation}\label{eq:rusanov}
\widehat{\mathcal F}_{\mathrm{Rus}}(U_L,U_R;\bm n)
=\tfrac12 A(\bm n)(U_L+U_R)-\tfrac12 a(U_R-U_L),\qquad a\geq1.
\end{equation}
由于本问题线性且 $A(\bm n)$ 对称，也可直接构造特征上风通量
\begin{equation}\label{eq:upwind}
\widehat{\mathcal F}_{\mathrm{up}}(U_L,U_R;\bm n)
=\tfrac12 A(\bm n)(U_L+U_R)
-\tfrac12\abs{A(\bm n)}(U_R-U_L),
\end{equation}
其中 $\abs A=R\,\operatorname{diag}(|\lambda_k|)R^\mathsf T$，$R$ 为正交特征向量矩阵。Rusanov 通量对零特征值方向也施加耗散，而精确特征上风通量不会。

\section{一维可复现的最小实验}
考虑周期区间 $x\in[0,L)$ 上的 \eqref{eq:1dwave}。令 $V=(E_y,B_z)^\mathsf T$，通量为 $f(V)=(B_z,E_y)^\mathsf T$。对宽度 $\Delta x$ 的均匀网格，用
\begin{equation}\label{eq:1dfv}
V_i^{n+1}=V_i^n-\frac{\Delta t}{\Delta x}
\bigl(\widehat f_{i+1/2}^n-\widehat f_{i-1/2}^n\bigr),
\end{equation}
其 Rusanov 通量在 $a=1$ 时为
\begin{equation}
\widehat f_{i+1/2}=\tfrac12\bigl(f(V_i)+f(V_{i+1})\bigr)-\tfrac12(V_{i+1}-V_i).
\end{equation}
初值取 $E_y(x,0)=B_z(x,0)=\exp[-(x-x_0)^2/\sigma^2]$（若采用周期边界，应使用其周期延拓），解析解为
\begin{equation}
E_y(x,t)=B_z(x,t)=\exp[-(x-t-x_0)^2/\sigma^2]
\quad\text{（在实线上，或取周期延拓）}.
\end{equation}
以 $L^1$ 或 $L^2$ 误差比较不同网格，并检查 CFL 数 $\nu=\Delta t/\Delta x$ 的影响。一阶显式格式取 $\nu\leq1$ 是此例的稳定性条件；多维、高阶或其他时间离散的稳定步长由相应算子和时间格式重新确定。

\begin{lstlisting}[caption={一维周期网格的 Rusanov 有限体积更新（伪代码）}]
Given cell averages V[i] = (Ey[i], Bz[i]), i = 0,...,N-1
Choose dx = L/N, dt with dt/dx <= 1
For each time step:
    For every interface i+1/2 (periodic indexing):
        VL = V[i];  VR = V[(i+1) mod N]
        Fhalf[i] = 0.5*(flux(VL)+flux(VR)) - 0.5*(VR-VL)
    For every cell i:
        Vnew[i] = V[i] - (dt/dx)*(Fhalf[i]-Fhalf[(i-1) mod N])
    V = Vnew
where flux(Ey,Bz) = (Bz,Ey)
\end{lstlisting}

\section{通量守恒与离散散度约束}
对六分量 $U$ 采用一般单元中心有限体积格式，可获得控制体通量平衡；离散 $\Div_h\bm B=0$ 还要求旋度与散度的相容关系。二者对应如下结构：
\begin{center}
\begin{tabular}{p{0.41\linewidth}p{0.43\linewidth}}
\toprule
通量保守性 & 散度约束保持\\
\midrule
相邻网格共用面通量，内部通量抵消 & 离散散度与离散旋度满足 $D_hC_h=0$\\
主要来自控制体积分与通量配对 & 主要来自拓扑关联关系（边--面--体）\\
\bottomrule
\end{tabular}
\end{center}
Maxwell 求解器应分别验证面通量平衡与散度约束传播。

\section{Yee 网格与 constrained transport 的几何核心}
\begin{askai}
在一个取向明确的立方单元上，将六个面磁通量按外法向求和，再用边电场线积分代入各面更新式。请逐条列出每条边的两次贡献，验证磁通量总和的时间导数为零。

进一步展开：若相邻面使用的同一条边电场不一致，抵消在哪一步失败？
\end{askai}
直角交错网格上，把磁通量 $B_f$ 放在网格\emph{面}，把电场环量 $E_e$ 放在网格\emph{边}。Faraday 方程在每一个有向面上的离散积分形式为
\begin{equation}\label{eq:CT}
\frac{\dd B_f}{\dd t}=-\sum_{e\subset\partial f}\operatorname{sgn}(f,e)E_e.
\end{equation}
这里 $B_f$ 表示穿过面的磁通量（不是点值），$E_e$ 表示沿边的电场线积分；$\operatorname{sgn}(f,e)$ 由取向决定。对某个网格体，把所有外表面的磁通量按有向求和：每条内部边在\emph{边界的边界}中出现两次且符号相反，所以
\begin{equation}\label{eq:dd0discrete}
\frac{\dd}{\dd t}\sum_{f\subset\partial K}\operatorname{sgn}(K,f)B_f=0.
\end{equation}
这就是格点上的 $D_hC_h=0$，因而在一致取向和边场共享条件下，磁散度约束保持到舍入误差量级。

\begin{remark}
Yee-FDTD 同时采用时间与空间交错；上述推导展示其中的\emph{空间离散拓扑}及磁约束保持机制。电 Gauss 约束需要兼容的电流沉积或离散连续性方程，完整实现还需处理边界、材料界面、稳定性和起始半步数据。
\end{remark}

\section{散度清理（divergence cleaning）}
如果采用单元中心有限体积格式，也可加入辅助变量，把散度误差转化为可传播、可耗散的模式（例如 GLM 清理）。这种方法适应已有通量求解框架，但会修改演化系统；它与严格利用离散复形保持 $D_hC_h=0$ 的 constrained transport 思想不同。这里略去辅助方程、清理波速和边界闭合；可检索 hyperbolic divergence cleaning、GLM Maxwell system，并要求给出完整能量与边界分析。

\section{建议的数值实践路线}
\begin{enumerate}
\item \textbf{一维：}实现 \eqref{eq:1dfv}，验证行波、波速、CFL 和网格收敛。
\item \textbf{二维：}扩展到 $U\in\RR^6$ 与法向数值通量，测量 $\norm{\Div_h\bm B}$ 随时间的变化。
\item \textbf{结构保持：}用边--面交错离散更新磁通，比较其离散散度与单元中心格式。
\item \textbf{带源：}单独验证离散连续性方程与电 Gauss 约束的相容性。
\end{enumerate}


\chapter{弱形式之后：Sobolev 视角、Galerkin 与高阶 DG 色散}\label{ch:dg}
\begin{chaptersummary}
上一章以有限体积展示通量与约束，本章从弱形式和函数空间出发构造 Galerkin 与 DG 离散。局部多项式和数值迹产生半离散系统，其能量与波传播误差还取决于通量和时间推进。通过一维可复算例子，本章将空间阶数、频散与计算代价联系起来。
\end{chaptersummary}
\studyroute{这里略去边元基函数与交换投影的构造，集中于 DG 弱形式、能量和一维频散。}{N\'ed\'elec elements; commuting projection; DG energy estimate; Fourier--Bloch symbol}{从最低阶边环量自由度解释切向连续性，再对 $p=1$ 的 DG 单元符号求物理分支。}
\section{从弱形式到近似空间：数学问题变成线性代数}
第 5 章的弱形式 \eqref{eq:weakinternal} 已经允许解不再处处可微。
本章据此选择函数空间与有限维近似。完备性、迹定理和嵌入定理的证明在此略去，可按 Sobolev trace theorem 与 Galerkin approximation 检索。
以简单的区域 $\Omega$ 为例，$L^2(\Omega)$ 提供平方可积的能量度量，
$H^1(\Omega)$ 要求一阶弱导数在 $L^2$ 中。
Maxwell 的旋度结构更适合 $H(\curl;\Omega)$，电/磁通量的法向迹
则自然关联 $H(\Div;\Omega)$；这两个空间通常不能直接由 $H^1$ 替代。

取空间网格 $\mathcal T_h=\{K\}$。
\emph{连续 Galerkin/FEM} 要求相邻单元在特定迹意义下相容：
例如 N\'ed\'elec 边元用于 $H(\curl)$ 保切向连续，Raviart--Thomas 元用于
$H(\Div)$ 保法向连续。其意义不仅在于插值精度，更在于与
\begin{equation}\label{eq:derhamfem}
H^1\xrightarrow{\grad}H(\curl)
\xrightarrow{\curl}H(\Div)\xrightarrow{\Div}L^2
\end{equation}
的微分复形结构相兼容；在合适离散空间上也希望相应算子复形成立。
\emph{不连续 Galerkin}（discontinuous Galerkin, DG）采取另一种选择：
每个单元都独立保存高阶多项式，跨单元耦合通过\emph{单值数值通量}实现。

\section{先看 Maxwell 的旋度型弱形式}\label{sec:maxwellweak}
\begin{studentinput}
设介质界面两侧 $\epsilon^-\ne\epsilon^+$，无自由面电荷，因此法向电位移连续。若额外要求电场各分量连续，会对法向电场施加什么限制？这如何帮助理解 $H^1$ 与 $H(\curl)$ 的切向连续性要求？

\textbf{你的答案：}\vspace{2cm}
\end{studentinput}

\begin{engineernote}
兼容离散按方程所需的迹选择自由度：边上的环量对应切向信息，面上的通量对应法向信息。标准 nodal DG 使用分片空间，跨单元连续性由数值迹或通量处理；通量因此同时影响界面耦合与离散结构。
\end{engineernote}
对真空、常系数且光滑的情形（允许给定电流 $\bm j$），取
$\bm v,\bm w\in H(\curl;\Omega)$ 为测试场，将
$\partial_t\bm E=\curl\bm B-\bm j$ 和
$\partial_t\bm B=-\curl\bm E$ 分别乘测试场并积分。
利用旋度分部积分恒等式
\begin{equation}
\int_\Omega(\curl\bm a)\cdot\bm v\,\dd x
=\int_\Omega\bm a\cdot\curl\bm v\,\dd x
+\int_{\partial\Omega}(\bm n\times\bm a)\cdot\bm v\,\dd S,
\end{equation}
得到
\begin{align}\label{eq:maxwellweakcurl}
(\partial_t\bm E,\bm v)_\Omega
-(\bm B,\curl\bm v)_\Omega
-\langle\bm n\times\bm B,\bm v\rangle_{\partial\Omega}
&=-(\bm j,\bm v)_\Omega,\\
(\partial_t\bm B,\bm w)_\Omega
+(\bm E,\curl\bm w)_\Omega
+\langle\bm n\times\bm E,\bm w\rangle_{\partial\Omega}
&=0.
\end{align}
光滑情形下，这就是对原 PDE 的一次分部积分改写；对低正则性场，
边界配对由\emph{切向迹及其对偶空间}解释。
实施 DG 时，每个单元分别使用这两式，并将内侧/外侧的切向场通过
\emph{数值切向通量}联接；这正是前一章“界面通量”在电磁场上的版本。
为让读者可以亲手推导矩阵和频散关系，下面先研究与 Maxwell 横向特征场
完全等价的一维对称双曲系统。

\section{分片 Sobolev 空间与局部弱形式}\label{sec:broken}
以一维为入口，定义分片空间
\begin{equation}
H^1(\mathcal T_h)=\{v:\;v|_K\in H^1(K),\ K\in\mathcal T_h\},
\quad
V_h^p=\{v:\;v|_K\in\mathbb P_p(K)\},
\end{equation}
其中 $\mathbb P_p(K)$ 是次数不高于 $p$ 的多项式。
在单元界面上，左、右迹 $v^-,v^+$ 可以不同，定义
$\{v\}=(v^-+v^+)/2$ 与 $[v]=v^+-v^-$（此处是一维取向约定）。
一般多维网格的外法线随单元变化，不可机械照搬标量跳跃的符号。

考虑一维常系数对称系统
\begin{equation}\label{eq:dgmodel}
\partial_tU+A\partial_xU=0,\qquad A=A^\mathsf T,
\end{equation}
它既是简化模型，也是第 4 章中 Maxwell 沿固定方向的法向系统。
设 $K_j$ 是一个一维网格单元，其两端点分别为
$x_{j-1/2}$ 和 $x_{j+1/2}$。取 $v_h\in(V_h^p)^m$，分部积分得：
\begin{equation}\label{eq:dgweak}
\int_{K_j}(\partial_tU_h)\cdot v_h\,\dd x
-\int_{K_j}(AU_h)\cdot\partial_xv_h\,\dd x
+\widehat f_{j+1/2}\cdot v_h(x_{j+1/2}^-)
-\widehat f_{j-1/2}\cdot v_h(x_{j-1/2}^+)=0.
\end{equation}
此处法向方向的正号约定统一向右；同一内部界面的
$\widehat f$ 对左右两个单元取\emph{相同}数值。
选一组单元局部基函数 $\{\phi_\ell\}_{\ell=0}^p$，展开
$U_h|_{K_j}=\sum_\ell u_{j\ell}(t)\phi_\ell(x)$，
即可得到局部质量矩阵和刚度矩阵，进而组成
$M_h\dot u_h=R_h(u_h)$。
使用局部正交 Legendre 多项式，质量矩阵在常系数、仿射一维单元上是对角的。

\section{中心通量、上风通量与离散能量}\label{sec:dg-energy}
\begin{askai}
对周期一维系统 $\partial_tU+A\partial_xU=0$，取对称常矩阵 $A$，从式 \eqref{eq:dgweak} 开始分别代入中心和上风通量，逐界面推导式 \eqref{eq:dgenergy}。请严格沿用 $[U]=U^+-U^-$ 的约定。

进一步展开：为什么半离散能量稳定仍需分别检查全离散时间推进与 Maxwell 散度约束？
\end{askai}
对 \eqref{eq:dgmodel}，两个常见的界面通量是
\begin{align}\label{eq:dg-fluxes}
\widehat f_{\rm cen}&=A\{U_h\},\\
\widehat f_{\rm up}&=A\{U_h\}-\frac12|A|[U_h],
\end{align}
其中 $|A|=R\operatorname{diag}(|\lambda_\alpha|)R^\mathsf T$，
$A=R\operatorname{diag}(\lambda_\alpha)R^\mathsf T$。
以周期边界、精确积分、常对称矩阵 $A$ 为假设，
在 \eqref{eq:dgweak} 中令 $v_h=U_h$，对所有单元求和。
单元体积项满足
$\int_K (AU_h)\cdot\partial_xU_h\,\dd x
=\tfrac12[U_h^\mathsf TAU_h]_{\partial K}$；
界面求和后得到
\begin{equation}\label{eq:dgenergy}
\frac{\dd}{\dd t}\left(\frac12\sum_K\norm{U_h}_{L^2(K)}^2\right)
=-\frac12\sum_{f\in\mathcal F_h^{\rm int}}
[U_h]_f^\mathsf T|A|[U_h]_f\leq0
\quad\text{（上风通量）}.
\end{equation}
中心通量时右端为零。高阶 DG 的稳定性取决于通量、积分精度、边界与时间推进。
上风通量引入的能量耗散通常用于压制不需要的网格尺度模式。
对于多维 Maxwell，能量估计仍可使用每个面对应的法向 $A(n)$，
但物理边界通量必须一并计入，而且散度约束仍需单独检查。

\section{高阶 DG 的数值频散关系}\label{sec:dg-dispersion}
\begin{checkpoint}
请比较振幅耗散、相速度误差和群速度误差；为什么中心通量可能守恒半离散能量，却仍有非零相位误差？能量估计为何也不能单独保证 Maxwell 离散谱正确？
\end{checkpoint}
数值波的传播精度需要同时比较振幅与相位。
设连续行波 $u(x,t)=\exp[\mathrm i(kx-\omega t)]$，真空中的一维特征波满足
$\omega=ck$。若某空间离散的物理模态满足
\begin{equation}\label{eq:semidispersion}
\lambda_h(k)=-\alpha_h(k)-\mathrm i\omega_h(k),
\quad
c_{\rm ph,h}=\frac{\omega_h}{k},\quad
c_{\rm gr,h}=\frac{\dd\omega_h}{\dd k},
\end{equation}
则 $\alpha_h\geq0$ 描述\emph{数值耗散}，
$\omega_h/k-c$ 描述\emph{数值相速度偏差（色散）}，
$\dd\omega_h/\dd k-c$ 描述\emph{群速度偏差}。
复杂电磁脉冲包含许多波数，若这些波数被赋予错误的相速度/群速度，
长期传播就会出现非物理波形展宽、相位漂移。

\Needspace{9\baselineskip}
\noindent\textbf{可复算的 Fourier--Bloch 单元符号。}\par
\begin{askai}
取一维周期均匀网格和 $p=1$ Legendre 基函数，从弱形式构造 $M,N,r,\ell$，再推导式 \eqref{eq:p1symbol}。求其两个特征分支，并解释如何在 $\theta\to0$ 时识别物理分支。

进一步展开：用附录代码复算 $\theta=0.5,1.0,1.5$ 的相速度误差，并与表格比较。
\end{askai}
取标量右行特征场
$\partial_tu+\partial_xu=0$、等距周期单元 $h$，
在参考区间 $\xi\in[-1,1]$ 上以 Legendre 多项式 $P_\ell(\xi)$
展开 $u_h$，采用\emph{精确积分和上风通量}。令
$\theta=kh$，$z=e^{-\mathrm i\theta}$，
单元系数采用 Bloch 形式
$u_{j,\ell}(t)=\widehat u_\ell(t)e^{\mathrm ij\theta}$。
设 $i,j=0,\ldots,p$，定义
\begin{equation}\label{eq:dgsymbol}
M_{ij}=\frac2{2i+1}\delta_{ij},\qquad
N_{ij}=\int_{-1}^1P_j(\xi)P_i'(\xi)\,\dd\xi,
\qquad
r_j=1,\quad\ell_i=(-1)^i.
\end{equation}
弱形式的系数向量满足
\begin{equation}\label{eq:dgsymboleig}
\frac{\dd\widehat{\bm u}}{\dd t}
=\frac1h\mathcal L_p(\theta)\widehat{\bm u},\qquad
\boxed{\mathcal L_p(\theta)=2M^{-1}\left(N-r r^\mathsf T
+z\,\ell r^\mathsf T\right).}
\end{equation}
该矩阵共有 $p+1$ 个特征分支：其中在 $\theta\to0$ 时连续趋于 $-\mathrm i\theta$
的分支称为\emph{物理分支}，其余是单元内部的非物理/高频分支；
绘制色散曲线需沿波数连续跟踪所选分支。

\paragraph{两个极端例子。}
当 $p=0$ 时，
\begin{equation}
\mathcal L_0(\theta)=e^{-\mathrm i\theta}-1,
\qquad
\omega_h=\frac{\sin\theta}{h},\quad
\alpha_h=\frac{1-\cos\theta}{h}.
\end{equation}
因此相速度比值为 $c_{\rm ph,h}/c=\sin\theta/\theta$，
随波数增大明显偏离 1。
当 $p=1$ 时，代入 Legendre 基函数得到
\begin{equation}\label{eq:p1symbol}
\mathcal L_1(\theta)=
\begin{pmatrix}z-1&z-1\\3(1-z)&-3(1+z)\end{pmatrix}.
\end{equation}
分别求取这两个矩阵以及 $p=2$ 矩阵的物理特征值，
可获得下表：每个数字都是 $|\omega_h/k-1|$，
只包含\emph{半离散空间误差}，取 $h=1$，无量纲波数为 $\theta=kh$。

\begin{center}
\begin{tabular}{cccc}
\toprule
$\theta=kh$ & $p=0$ & $p=1$ & $p=2$\\
\midrule
$0.5$ & $4.11489\times10^{-2}$ & $2.17952\times10^{-4}$ & $3.68086\times10^{-7}$\\
$1.0$ & $1.58529\times10^{-1}$ & $2.96135\times10^{-3}$ & $2.27506\times10^{-5}$\\
$1.5$ & $3.35003\times10^{-1}$ & $1.18404\times10^{-2}$ & $2.42693\times10^{-4}$\\
\bottomrule
\end{tabular}
\end{center}

\begin{figure}[htbp]
\centering
\begin{tikzpicture}
\begin{semilogyaxis}[
 width=0.90\textwidth,height=6.6cm,
 xmin=0.2,xmax=1.55,ymin=1e-9,ymax=1,
 xlabel={$\theta=kh$},ylabel={相速度相对误差 $|c_{\mathrm{ph},h}/c-1|$},
 grid=major,grid style={gray!25},
 legend pos=north west,
 tick label style={font=\footnotesize},
 label style={font=\small},
 legend style={font=\small,fill=white,draw=gray!45},
 ]
\addplot[blue!70!black,thick,mark=square*] coordinates {
 (0.25,0.010384163) (0.5,0.041148923) (0.75,0.09114832)
 (1,0.158529015) (1.25,0.240812305) (1.5,0.335003342)};
\addlegendentry{$p=0$}
\addplot[orange!80!black,thick,mark=triangle*] coordinates {
 (0.25,1.42467155e-5) (0.5,0.000217952365) (0.75,0.00102789518)
 (1,0.00296135335) (1.25,0.00647412633) (1.5,0.0118404387)};
\addlegendentry{$p=1$}
\addplot[teal!80!black,thick,mark=*] coordinates {
 (0.25,5.79772053e-9) (0.5,3.68085612e-7) (0.75,4.13434552e-6)
 (1,2.27506265e-5) (1.25,8.43451092e-5) (1.5,0.000242693401)};
\addlegendentry{$p=2$}
\end{semilogyaxis}
\end{tikzpicture}
\caption{一维常系数上风 DG 的半离散物理模态色散误差（纵坐标对数刻度）。
数据按式 \eqref{eq:dgsymboleig} 的特征值逐点计算，适用于本节指定网格、积分与通量，仅包含空间半离散误差。}
\label{fig:dgdisp}
\end{figure}
\FloatBarrier

\noindent\textbf{误差趋势。}
在这一\emph{特定}均匀网格、常系数、精确积分、特征上风通量的例子中，
从常数近似提高到一次、二次局部多项式，\emph{物理模态}的相速度误差显著降低（见图~\ref{fig:dgdisp}）。
对小 $kh$，相对色散误差依次具有
$O((kh)^2)$、$O((kh)^4)$、$O((kh)^6)$ 的行为。
复杂网格、材料跳跃、欠积分以及其他通量会改变频散曲线，上述规律适用于本例所列条件。
另外，$p$ 提高会增加局部未知量、积分代价并收紧显式时间步长；
同等\emph{总自由度或计算成本}下的比较需要独立的基准实验。

\section{空间与时间误差的组合}
令物理空间特征值为 $\lambda_h=-\alpha_h-\mathrm i\omega_h$，
时间积分器的放大因子为 $R(\Delta t\lambda_h)$，则完整的离散频率是
\begin{equation}\label{eq:fullfrequencies}
\alpha_{h,\Delta t}=-\frac1{\Delta t}\log|R(\Delta t\lambda_h)|,
\qquad
\omega_{h,\Delta t}=-\frac1{\Delta t}\arg R(\Delta t\lambda_h),
\end{equation}
取连续的相位分支以避免 $2\pi$ 混叠。
无论 DG 空间阶数多高，如果时间格式很低阶或 $\Delta t$ 太大，
$\omega_{h,\Delta t}$ 仍可能偏离连续模型的频率。
中点法的式 \eqref{eq:midpointdisp} 给出了时间色散的具体表达。
这也解释了工程中的常见搭配：高阶 DG + 与其精度匹配的 RK 时间格式；
或在长期结构保持优先时，采用针对离散 Hamilton/Poisson 结构设计的推进策略。
\emph{数值色散小}、\emph{数值耗散小}、\emph{离散能量/辛结构好}是互相关联但不能互相替代的评价维度。

\section{三维 Maxwell DG 的空间与约束设计}
Maxwell DG 通常取每个单元的电场与磁场分量为局部多项式，
以 $A(n)$ 的中心或特征上风通量处理面耦合；
可以支持高阶曲面几何、非结构网格及 $hp$ 局部加密。
逐单元 $[\mathbb P_p]^6$ 空间还需配套约束设计，以控制磁散度和带源电 Gauss 约束。可采用离散复形兼容空间、
适当的约束保持通量、局部散度自由基、清理变量或其它专门设计。
物理边界和材料界面的稳定性应在相应通量下分析。

\section{高阶色散的数值实验}
\begin{enumerate}
\item 取 $p=0,1,2,3$，按 \eqref{eq:dgsymbol} 生成 $M,N,r,\ell$，取
$\theta=0.1,0.2,\ldots,1.5$；对 $\mathcal L_p(\theta)$ 求特征值，追踪
与 $-\mathrm i\theta$ 连续相接的物理模态。
\item 画 $c_{\rm ph,h}/c$、$\alpha_hh$ 随 $kh$ 的曲线；
比较不同 $p$ 的相误差，留意非物理分支不能误选。
\item 再选 RK4 或隐式中点法，用 \eqref{eq:fullfrequencies} 在多个
$\Delta t/h$ 下重新画\emph{全离散}色散曲线，判断何时由时间误差主导。
\item 回到三维高斯脉冲，比较长期传播后的相位、能量与 $\Div_h\bm B$。
如使用不同网格和多项式次数，需同时报告总自由度和实际运行成本。
\end{enumerate}
\section{练习}
\begin{enumerate}
\item 从 \eqref{eq:dgweak} 出发，证明 $p=0$ 的 DG--上风格式就是一阶上风有限体积格式。
\item 对 \eqref{eq:p1symbol} 求特征多项式并找出 $\theta\to0$ 的物理分支。
\item 由中心通量推导周期 DG 半离散能量严格守恒；解释加上上风耗散项后的变化。
\item 讨论在 $p=2$、固定空间网格下，仅把 RK4 换成显式 Euler 会怎样影响纯振荡模态的稳定性。
\end{enumerate}

\chapter*{结语：从几何到计算的闭环}
\addcontentsline{toc}{chapter}{结语：从几何到计算的闭环}
本文的逻辑可压缩为一条链：
\begin{equation*}
\boxed{
\begin{gathered}
\text{Lorentz 协变性与统一场强 }F\\
\Downarrow\\
\dd F=0,\qquad \dd\star F=J\\
\Downarrow\quad\text{选定时间分叶}\\
\text{6 个演化分量 + 2 个约束分量}\\
\Downarrow\quad\dd^2=0\\
\text{约束传播与源的连续性条件}\\
\Downarrow\\
\text{对称双曲系统 / 能量估计 / 守恒型}\\
\Downarrow\\
\text{数值通量 + 保结构时间推进 + 高阶 DG 与保约束离散}
\end{gathered}}
\end{equation*}

本卷的主线可概括为：\emph{几何结构决定方程的相容性；双曲结构决定信息如何传播；离散结构决定计算能否保留这些性质。} Maxwell 方程由此连接几何场论与计算 PDE。

\appendix
\chapter{符号与约定速查}
\begin{chaptersummary}
本附录汇总前文的时空、微分形式与离散算子记号，便于在四维表达、三维分解和数值实现之间切换。使用公式时应同时核对度规、取向和单位约定。
\end{chaptersummary}
\studyroute{这里汇总前文约定；一般材料本构与其他号差下的公式转换可单独展开。}{metric signature; electromagnetic constitutive relations; Hodge dual}{保持本卷取向，说明恢复 $\varepsilon,\mu,c$ 后能量内积与波速如何改变。}
\begin{longtable}{p{0.26\textwidth}p{0.66\textwidth}}
\toprule
记号 & 约定与解释\\
\midrule
\endfirsthead
\toprule
记号 & 约定与解释\\
\midrule
\endhead
$\eta$ & $(-,+,+,+)$ 的 Minkowski 度规，$c=1$。\\
$\Omega^k(M)$ & 四维时空上的 $k$ 形式；$\Omega^k(\Sigma)$ 为空间 $k$ 形式。\\
$\dd,\dd_\Sigma$ & 四维/空间外微分。\\
$\star,\starS$ & 四维 Lorentz / 三维欧氏 Hodge 星。\\
$E$ & 空间电场一形式：$E_i\dd x^i$。\\
$B$ & 空间磁通量二形式：$B_x\dd y\wedge\dd z+\cdots$。\\
$F$ & $F=B-\dd t\wedge E$，所以 $F_{0i}=-E_i$。\\
$J$ & 电流三形式：$J=\rho\vol_3-\dd t\wedge\starS j^\flat$。\\
$\star F$ & $\starS E+\dd t\wedge\starS B$。\\
$C(\bm n)$ & $C(\bm n)\bm v=\bm n\times\bm v$。\\
$A(\bm n)$ & 对称 $6\times6$ 主部矩阵，特征值 $\pm1,0$。\\
$D_hC_h=0$ & 离散散度作用在离散旋度上为零，依赖同一离散复形。\\
\bottomrule
\end{longtable}

\paragraph{三种不同的“星”或对偶}这里的 Hodge 星不等于复共轭，也不等于矩阵伴随符号 $A^*$。四维 Lorentz Hodge 星与三维欧氏 Hodge 星在平方后的符号一般不同；不要直接把 $\star^2=1$ 从三维照搬到四维。

\paragraph{度规与单位}把 $c$ 和材料本构恢复后，$\bm E$、$\bm B$ 之间还涉及物理单位与本构关系。对非均匀介质，最自然的变量可能是 $\bm D,\bm H$；主部矩阵、能量内积和特征速度都会随本构变化。本卷的 $\pm1$ 传播速度只适用于所选单位制下的真空平直时空。

\chapter{DG 数值频散复算：Python 代码}
\begin{chaptersummary}
本附录把第 7 章的一维 DG 单元符号转化为可运行的特征值计算，用于复现物理分支的半离散相速度误差。改变波数与多项式阶数，可以检验分析结论并观察分支选择的影响。
\end{chaptersummary}
\studyroute{这里提供单元符号的数值复算，略去完整 DG 场求解器的装配与时间推进。}{Legendre modal DG; physical eigenvalue branch; dispersion analysis}{复现表中误差，再沿波数连续跟踪特征值，比较最近特征值选择与连续跟踪。}
本附录代码只需要 NumPy；针对一维常系数、周期均匀网格、精确 Legendre 积分、上风通量
和\emph{空间半离散}。它构造式 \eqref{eq:dgsymboleig} 并输出第 7 章的相速度误差表。
代码分析与 Maxwell 横向特征场等价的标量行波；完整 Maxwell 场和 Gauss 约束不在这一实验范围内。

\Needspace{30\baselineskip}
\begin{lstlisting}[language=Python,caption={Legendre 模态 DG 的空间色散：物理特征值分析}]
import numpy as np
from numpy.polynomial.legendre import Legendre, leggauss

def physical_mode(p, theta):
    x, w = leggauss(max(10, 2*p + 4))
    B  = np.array([Legendre.basis(j)(x) for j in range(p+1)])
    dB = np.array([Legendre.basis(i).deriv()(x) for i in range(p+1)])
    N = (dB * w) @ B.T   # N[i,j] = integral P_j * P_i'
    ii = np.arange(p+1)
    mass = 2.0 / (2*ii + 1)
    r = np.ones(p+1)
    ell = (-1.0)**ii
    z = np.exp(-1j*theta)
    L = (2.0 / mass[:, None]) * (
        N - np.outer(r, r) + z*np.outer(ell, r))
    eigenvalues = np.linalg.eigvals(L)
    # Valid for this moderate theta range: nearest to exact -i*theta.
    return eigenvalues[np.argmin(abs(eigenvalues + 1j*theta))]

for theta in (0.5, 1.0, 1.5):
    errors = []
    for p in (0, 1, 2):
        lam = physical_mode(p, theta)
        error_phase = abs(-lam.imag/theta - 1.0)
        errors.append(error_phase)
    print(theta, *(f'{err:.6e}' for err in errors))
\end{lstlisting}
在 $\theta$ 接近周期 Brillouin 区边缘、物理模态与其它分支变得接近时，
应采用沿波数连续追踪特征分支的方法，而不是始终简单选择距离
$-\mathrm i\theta$ 最近的特征值。

\chapter{参考读物与扩展方向}
\begin{chaptersummary}
本附录按电磁学、微分形式、PDE、有限体积、几何积分与 DG 等方向列出延伸读物，供读者针对正文中需要深化的环节继续学习。比较公式时应留意不同文献采用的符号约定。
\end{chaptersummary}
\studyroute{以下文献按主题提供进一步阅读入口，具体证明可结合前文关键词定位。}{computational electromagnetism; geometric integration; discontinuous Galerkin}{按所需证明或算法列出阅读顺序，并将文献记号与本卷符号表逐项对应。}
以下读物按知识环节提供进一步查阅的入口；比较公式时需统一度规号差、电磁张量分量与 Hodge 星约定。
\begin{enumerate}
\raggedright
\item D. J. Griffiths, \emph{Introduction to Electrodynamics}. 适于复习三维 Maxwell 方程和电磁波。
\item B. F. Schutz, \emph{A First Course in General Relativity}. 适于学习四维张量和简单的时空几何。
\item S. A. Hughes, \href{https://ocw.mit.edu/courses/8-033-introduction-to-relativity-and-spacetime-physics-fall-2024/mit8_033_f24_lec_full.pdf}{\emph{Introduction to Relativity and Spacetime Physics}}, MIT 8.033 (2024), Lectures 11--12.
适于核对惯性系间的电磁场变换与 Maxwell 方程的张量表达。
\item H. Flanders, \emph{Differential Forms with Applications to the Physical Sciences}. 适于从向量分析走向外微分形式。
\item D. Tong, \href{https://www.davidtong.org/teaching/general-relativity/grhtml/S2}{\emph{General Relativity}, Chapter 2}.
第 2.3.1、2.3.2、2.4.1 节分别介绍对偶坐标基、形式的 Lie 导数与外微分。
\item M. Fecko, \href{https://arxiv.org/abs/gr-qc/9701066}{\emph{On 3+1 decompositions with respect to an observer field via differential forms}},
\emph{Journal of Mathematical Physics}, 38 (1997), 4542--4560,
\href{https://doi.org/10.1063/1.532142}{doi:10.1063/1.532142}.
第 IV.B 节用空间外微分与 Lie 导数给出时空分解；本文采用的固定乘积坐标对应其简化情形。
\item L. C. Evans, \emph{Partial Differential Equations}. 适于补充弱解、函数空间和能量方法。
\item R. J. LeVeque, \emph{Finite Volume Methods for Hyperbolic Problems}. 适于系统学习守恒律、数值通量和波传播。
\item A. Bossavit, \emph{Computational Electromagnetism: Variational Formulations, Complementarity, Edge Elements}. 适于理解电磁场计算中几何离散的作用。
\item E. Hairer, C. Lubich, G. Wanner, \emph{Geometric Numerical Integration: Structure-Preserving Algorithms for Ordinary Differential Equations}. 适于深入理解保辛、长时间能量行为与数值色散。
\item J. S. Hesthaven and T. Warburton, \emph{Nodal Discontinuous Galerkin Methods: Algorithms, Analysis, and Applications}. 适于实现 DG 半离散及高阶局部基函数。
\item P. Monk, \emph{Finite Element Methods for Maxwell's Equations}. 适于学习 $H(\curl)$ 与兼容有限元的 Maxwell 离散。
\item M. Ainsworth, ``Dispersive and dissipative behaviour of high order discontinuous Galerkin finite element methods,'' \emph{Journal of Computational Physics}, 198 (2004), 106--130. 适于专题阅读 DG 数值频散与耗散。
\item D. S\'arm\'any, M. A. Botchev, and J. J. W. van der Vegt, ``Dispersion and Dissipation Error in High-Order Runge--Kutta Discontinuous Galerkin Discretisations of the Maxwell Equations,'' \emph{Journal of Scientific Computing}, 33 (2007), 47--74. 适于比较空间高阶与时间推进的综合频散特性。
\item K. S. Yee, ``Numerical Solution of Initial Boundary Value Problems Involving Maxwell's Equations in Isotropic Media,'' \emph{IEEE Transactions on Antennas and Propagation}, 1966. 适于追溯交错电磁网格格式的原始思想。
\item M. Dumbser (2014), \href{https://doi.org/10.1016/j.cma.2014.07.019}{\emph{Arbitrary-Lagrangian-Eulerian ADER-WENO Finite Volume Schemes with Time-Accurate Local Time Stepping for Hyperbolic Conservation Laws}}.
\item B. Owren and M. Zennaro (1992), \href{https://doi.org/10.1137/0913084}{\emph{Derivation of Efficient, Continuous, Explicit Runge--Kutta Methods}}.
\item C. Altmann, G. Gassner, and C.-D. Munz (2010), \href{https://congress2.cimne.com/eccomas/proceedings/cfd2010/papers/01008.pdf}{\emph{An Explicit Space-Time Adaptive Discontinuous Galerkin Scheme}}, ECCOMAS CFD.
\item U. M. Ascher, S. J. Ruuth, and R. J. Spiteri (1997), \href{https://doi.org/10.1016/S0168-9274(97)00056-1}{\emph{Implicit-explicit Runge--Kutta methods for time-dependent partial differential equations}}, Applied Numerical Mathematics, 25, 151--167.
\end{enumerate}

\end{document}
